\documentclass[12pt,a4paper]{article}

\usepackage[utf8]{inputenc}
\usepackage[T1]{fontenc}
\usepackage[ngerman]{babel}
\usepackage{amsmath, amssymb, amsthm}
\usepackage{mathtools}
\usepackage{geometry}
\usepackage{hyperref}
\usepackage{cleveref}
\usepackage{enumitem}
\usepackage{algorithm}
\usepackage{algpseudocode}
\usepackage{booktabs}
\usepackage{array}
\usepackage{graphicx}
\usepackage{tikz}
\usetikzlibrary{arrows.meta, positioning, calc}
\usepackage{pgfplots}
\pgfplotsset{compat=1.18}
\usepackage{float}
\usepackage{caption}
\usepackage{subcaption}
\usepackage{lmodern}
\usepackage{microtype}
\usepackage{setspace}
\usepackage{fancyhdr}
\usepackage{enumitem}
\usepackage{titlesec}
\usepackage{cite}

\geometry{margin=2.5cm}

\setlist[itemize,1]{label=--}

% Theorem-Umgebungen
\theoremstyle{definition}
\newtheorem{definition}{Definition}[section]
\newtheorem{beispiel}{Beispiel}[section]

\theoremstyle{plain}
\newtheorem{satz}{Satz}[section]
\newtheorem{lemma}[satz]{Lemma}
\newtheorem{korollar}[satz]{Korollar}
\newtheorem{theorem}[satz]{Theorem}

\theoremstyle{remark}
\newtheorem{bemerkung}{Bemerkung}[section]

% Operatoren
\DeclareMathOperator{\E}{\mathbb{E}}
\DeclareMathOperator{\Var}{Var}
\DeclareMathOperator{\Cov}{Cov}
\DeclareMathOperator{\Prob}{\mathbb{P}}

\hyphenation{Wahr-schein-lich-keit Wahr-schein-lich-keits-ver-tei-lung
  Be-ding-te Be-ding-ten Er-eig-nis-rau-mes Er-geb-nis-rau-mes}

\hypersetup{
  colorlinks=true,
  linkcolor=blue!60!black,
  citecolor=green!50!black,
  urlcolor=cyan!60!black,
  pdftitle={Die Lehre vom Zufall in der Mathematik},
  pdfauthor={},
  pdfsubject={Wahrscheinlichkeitstheorie und randomisierte Algorithmen}
}


\title{\textbf{Die Lehre vom Zufall in der Mathematik} \\[-0.1em]
\large Vollständige Beschreibbarkeit von Experimenten,\\
	Wahrscheinlichkeitstheorie und randomisierte Algorithmen}

\author{Stephan Epp}
\date{8. Januar 2026}

\begin{document}

\maketitle
\tableofcontents
\newpage

% ===============================================================
\section{Einleitung}
% ===============================================================
Diese Arbeit behandelt die mathematischen Grundlagen der Lehre vom Zufall.
Ausgangspunkt ist das fundamentale Prinzip der vollständigen Beschreibbarkeit
von Experimenten: Ein Experiment ist genau dann vollständig beschrieben, wenn
die Summe der Wahrscheinlichkeiten aller möglichen Ereignisse gleich $1$ ist –
sowohl im diskreten als auch im kontinuierlichen Wahrscheinlichkeitsraum.
Aufbauend auf diesem Grundprinzip werden Erwartungswert, Varianz,
bedingte Wahrscheinlichkeit, bedingter Erwartungswert und die wichtigsten
Wahrscheinlichkeitsverteilungen systematisch entwickelt. Im algorithmischen
Teil werden randomisierte Algorithmen der Informatik vorgestellt und tiefgehend
analysiert: der probabilistische $\frac{7}{8}$-Approximationsalgorithmus für
3-SAT, der Algorithmus von Karger für minimale Schnitte, randomisiertes
Quicksort, der Miller-Rabin-Primzahltest sowie der Algorithmenentwurf mittels
der Probabilistischen Methode. Abschließend wird ein umfangreicher Ausblick
auf offene Fragen und aktuelle Forschungsrichtungen gegeben.

Die Frage, was Zufall bedeutet, begleitet die Mathematik seit der Antike.
Bereits die griechischen Philosophen diskutierten, ob Ereignisse kausal
determiniert oder prinzipiell unvorhersehbar sein können. Pascal und Fermat
legten im 17. Jahrhundert mit ihrer Korrespondenz über das \emph{Prob\-lème
des points} den Grundstein für eine systematische mathematische Beschreibung
zufälliger Ereignisse \cite{david1962games}. Jacob Bernoulli formulierte das
Gesetz der großen Zahlen, Abraham de Moivre erkannte die Normalverteilung,
und Pierre-Simon Laplace schuf in seiner \emph{Théorie analytique des
probabilités} \cite{laplace1812theorie} eine erste axiomatische Fundierung.
Den entscheidenden Formalisierungsschritt vollzog Andrei Nikolajewitsch
Kolmogorow 1933 mit seiner axiomatischen Wahrscheinlichkeitstheorie
\cite{kolmogorov1933grundbegriffe}, die bis heute das Fundament des
Fachs bildet.

Wir verwenden Zufall, weil diese Art der Beschreibung weniger kompliziert ist als eine realistischere Beschreibung. Eine realistischere Beschreibung berücksichtigt weit mehr Randbedingungen und ist daher komplizierter und länger. Der Leser muss auch Aufwand zum Verstehen aufbringen. Der Wert zur Wiederverwendung ist dabei allerdings viel geringer, denn die Beschreibung durch Zufall hat viel mehr Anwendungen - vor allem im Entwurf von Algorithmen.

Gibt es einen Beweis in der Lehre vom Zufall, der mit Hilfe von Abbildungen bzw. Funktionen möglich ist, aber ohne nicht? Randbemerkung: Abbildungen eignen sich für diskrete Beschreibungen und Funktionen eignen sich für kontinuierliche Beschreibungen. 

Es ist sinnvoll, von der Modellierung in der Lehre vom Zufall zu sprechen. Denn Experimente werden in gewisser Weise, ob diskret oder kontinuierlich, abstrakt, also mit Hilfe eines Modells, beschrieben. Modelle abstrahieren und erlauben eine einfachere und kürzere Beschreibung. Modelle eignen sich besonders für die Wiederverwendung. Einen besonderen Nutzen haben sie dabei in der Informatik, wo sie auch für die Codegenerierung verwendet werden. 

Im Allgemeinen ist die Beschreibung durch ein Modell viel leichter in Code zu verfassen, das tut der Mensch oder ein KI-Chat. Dieser Code kann auf viele Instanzen oder Eingaben angewendet und automatisiert ausgeführt werden. Das ermöglicht das Ausführen von vielen Wiederholungen und unterschiedlichen Eingaben für allgemeinere und sicherere Aussagen in der konkreten Beschreibung oder Modellierung.

Was beschreibbar ist, ist berechenbar. Diese Aussage wurde soeben in der praktischen Umsetzung beschrieben. Dieses Vorgehen ist auf beliebige Beschreibungen oder Modelle anwendbar und der große Nutzen der Modellierung und Berechenbarkeit wird dadurch sehr deutlich.

Der zentrale Gedanke dieser Arbeit geht auf Kolmogorows drittes Axiom
zurück: \emph{Die Summe der Wahrscheinlichkeiten aller Elementarereignisse
eines Experiments ist 1.} Dieses scheinbar einfache Axiom hat weitreichende
Konsequenzen: Es legt fest, was ein \emph{vollständig beschriebenes Experiment}
ist, und es bildet die operationale Grundlage für alle weitergehenden Begriffe
wie Erwartungswert, Varianz und bedingte Wahrscheinlichkeit.

% ===============================================================
\section{Grundlagen}
\label{sec:grundlagen}
% ===============================================================

Um jede Unruhe über den Ausgang des Zufalls allen Beteiligten zu nehmen, wird
der Zufall für diskrete Ereignisse mit Hilfe der Mathematik beschrieben und
analysiert.

\subsection{Ergebnisraum und Ereignisse}

Für das diskrete Ereignis, bei dem ein Würfel einmal geworfen wird, können
folgende diskrete Ereignisse zufällig eintreten. Der Würfel kann auf die
Augenzahl $1$ fallen, der Würfel kann auf die Augenzahl $2$ fallen, \ldots,
der Würfel kann auf die Augenzahl $6$ fallen. Dabei ist jeder Fall des Würfels
auf eine andere Augenzahl ein Ereignis. Es gibt damit also sechs Ereignisse
beim Würfeln des Würfels mit sechs Seiten.

Ein \textit{Versuch} beschreibt das Vorhaben, zum Beispiel einen Würfel zu
würfeln. In diesem Versuch können unterschiedliche Ereignisse zufällig
eintreten. \textit{Die Menge aller Ereignisse für einen Versuch}, die
zufällig eintreten können, wird definiert mit $\Omega$. Um nun zu sagen, wie
zufällig es ist, dass der Würfel beim Wurf auf die Augenzahl $3$ fällt, muss
nur dieses Ereignis ins Verhältnis zu allen möglichen Ereignissen gesetzt
werden. Damit ergibt sich als Wert für den Zufall des Ereignisses, dass der
Würfel beim Wurf auf die Augenzahl $3$ fällt, die \textit{Wahrscheinlichkeit}
von $\frac{1}{6}$. Es fällt auf, dass für jedes Ereignis dieses Versuches
gilt, dass der Würfel beim Wurf auf die Augenzahl $k$ fällt, die
\textit{Wahrscheinlichkeit} $\frac{1}{6}$ hat, $k \in \Omega = \{1, \ldots, 6\}$.
Denn der Würfel hat ja sechs gleich große Seiten und keine der Seiten des
Würfels ist durch die Art des Würfels bevorzugt. Ist die Summe aller
Wahrscheinlichkeiten der Zufälle aller Ereignisse insgesamt $1$, wurden alle
Ereignisse des Versuchs berücksichtigt und der Versuch ist
\textit{mathematisch vollständig beschrieben}.

Für Versuche, die durch diskrete Ereignisse beschrieben werden, ist die
Überprüfung der vollständigen mathematischen Beschreibung notwendig. Erst dann
liegt eine mathematische Beschreibung vor, auf deren Grundlage weitere
Ereignisse für diesen Versuch betrachtet und analysiert werden können. Für
den Versuch, einen Würfel einmal zu würfeln, kann man sich auch fragen, wie
wahrscheinlich der Zufall für das Ereignis ist, dass der Würfel nicht auf die
Augenzahl $3$ fällt. Das ist komplizierter im Allgemeinen. Glücklicherweise
aber betrachten wir diskrete Ereignisse. Deshalb lässt sich die
Wahrscheinlichkeit für den Zufall des Ereignisses, dass der Würfel nicht auf
die Augenzahl $3$ fällt, wie folgt \textit{abzählen}. Die Wahrscheinlichkeit
für den Zufall des Ereignisses, dass der Würfel nicht auf die Augenzahl $3$
fällt, ist dieselbe wie für den Zufall der Ereignisse, dass der Würfel auf
die Augenzahl $1, 2, 4, 5$ oder $6$ fällt. Dies sind fünf Ereignisse. Damit
ergibt sich eine Wahrscheinlichkeit von $\frac{5}{6}$.

\begin{definition}[Ergebnisraum und Ereignis]\label{def:ergebnisraum}
  Sei $\Omega$ eine nichtleere Menge, die sog.\ \textbf{Ergebnismenge} oder
  \textbf{Stichprobenraum}. Ihre Elemente $\omega \in \Omega$ heißen
  \textbf{Elementarereignisse} oder \textbf{Ergebnisse}. Eine Teilmenge
  $A \subseteq \Omega$ heißt \textbf{Ereignis}. Die Potenzmenge
  $\mathcal{P}(\Omega)$ ist im diskreten Fall der \textbf{Ereignisraum};
  im allgemeinen Fall tritt an ihre Stelle eine $\sigma$-Algebra
  $\mathcal{F} \subseteq \mathcal{P}(\Omega)$.
\end{definition}

Ein Ereignis $A$ \emph{tritt ein}, wenn das realisierte Ergebnis
$\omega \in A$ liegt. Das sichere Ereignis ist $\Omega$ selbst, das
unmögliche Ereignis ist $\emptyset$.

\begin{definition}[$\sigma$-Algebra]\label{def:sigma-algebra}
  Eine Mengenfamilie $\mathcal{F} \subseteq \mathcal{P}(\Omega)$ heißt
  \textbf{$\sigma$-Algebra}, wenn gilt:
  \begin{enumerate}[label=(\roman*)]
    \item $\Omega \in \mathcal{F}$,
    \item $A \in \mathcal{F} \Rightarrow A^c \coloneqq \Omega \setminus A \in \mathcal{F}$,
    \item $A_1, A_2, \ldots \in \mathcal{F} \Rightarrow \bigcup_{i=1}^{\infty} A_i \in \mathcal{F}$.
  \end{enumerate}
  Das Tripel $(\Omega, \mathcal{F})$ heißt \textbf{Messbarer Raum}.
\end{definition}

\subsection{Kolmogorows Axiome und das Prinzip der Vollständigkeit}
\label{subsec:axiome}

\begin{definition}[Wahrscheinlichkeitsmaß (Kolmogorow, 1933)]\label{def:wmaß}
  Sei $(\Omega, \mathcal{F})$ ein messbarer Raum. Eine Abbildung
  $\Prob\colon \mathcal{F} \to [0,1]$ heißt \textbf{Wahrscheinlichkeitsmaß},
  wenn die drei Kolmogorowschen Axiome gelten:
  \begin{enumerate}[label=\textbf{(K\arabic*)}]
    \item \textbf{Nichtnegativität:} $\Prob(A) \geq 0$ für alle $A \in \mathcal{F}$.
    \item \textbf{Normierung:} $\Prob(\Omega) = 1$.
    \item \textbf{$\sigma$-Additivität:} Für paarweise disjunkte Ereignisse
      $A_1, A_2, \ldots \in \mathcal{F}$ gilt
      \[
        \Prob\!\left(\bigcup_{i=1}^{\infty} A_i\right) = \sum_{i=1}^{\infty} \Prob(A_i).
      \]
  \end{enumerate}
  Das Tripel $(\Omega, \mathcal{F}, \Prob)$ heißt \textbf{Wahrscheinlichkeitsraum}.
\end{definition}

Das zweite Axiom (K2) – die \emph{Normierung} – ist das mathematische Fundament
des folgenden zentralen Prinzips dieser Arbeit:

\paragraph*{Prinzip der vollständigen Beschreibbarkeit.}
\begin{quote}
  \textbf{Diskreter Fall.} Ein Experiment mit diskretem Ergebnisraum
  $\Omega = \{\omega_1, \omega_2, \ldots\}$ ist \emph{genau dann vollständig
  beschrieben}, wenn
  \[
    \sum_{i} \Prob(\{\omega_i\}) = 1.
  \]
  \textbf{Kontinuierlicher Fall.} Ein Experiment mit stetigem Ergebnisraum
  $\Omega \subseteq \mathbb{R}^n$ und Dichtefunktion $f\colon \Omega \to [0,\infty)$
  ist \emph{genau dann vollständig beschrieben}, wenn
  \[
    \int_{\Omega} f(\omega)\, d\omega = 1.
  \]
  In beiden Fällen besagt die vollständige Beschreibbarkeit: Jedes mögliche
  Ergebnis des Experiments ist berücksichtigt; kein Ergebnis ist \glqq
  vergessen\grqq{} worden, und kein Szenario wird doppelt gezählt.
\end{quote}

\begin{beispiel}[Würfelwurf – vollständige Beschreibung]\label{bsp:wuerfel}
  Betrachte einen fairen sechsseitigen Würfel. Der Ergebnisraum ist
  $\Omega = \{1, 2, 3, 4, 5, 6\}$. Das Experiment \emph{„Wirf den Würfel
  einmal"} weist jedem Elementarereignis die gleiche Wahrscheinlichkeit zu:
  \[
    \Prob(\{k\}) = \frac{1}{6}, \quad k \in \{1,\ldots,6\}.
  \]
  Die Vollständigkeitsbedingung ist erfüllt:
  \[
    \sum_{k=1}^{6} \Prob(\{k\}) = 6 \cdot \frac{1}{6} = 1.
  \]
  Das Teilexperiment \emph{„Der Würfel zeigt die } 1 hat den Ergebnisraum
  $\Omega' = \{1\}$ mit $\Prob(\{1\}) = \frac{1}{6}$. Diese Beschreibung ist
  \emph{unvollständig}, denn $\frac{1}{6} \neq 1$. Um sie zu vervollständigen,
  fügen wir das Komplementärereignis $\{2,3,4,5,6\}$ hinzu:
  \[
    \Prob(\{1\}) + \Prob(\{2,3,4,5,6\}) = \frac{1}{6} + \frac{5}{6} = 1.
  \]
  Erst durch das explizite Benennen des Komplementärereignisses – \emph{„der
  Würfel zeigt \textbf{nicht} die }1 – ist das Experiment vollständig
  beschrieben. Mit Wahrscheinlichkeit $\frac{1}{6}$ fällt der Würfel auf die
  Augenzahl~$1$ und mit Wahrscheinlichkeit $\frac{5}{6}$ fällt der Würfel
  nicht auf die Augenzahl~$1$.
\end{beispiel}

\subsection{Wichtige Rechenregeln}

Aus den Kolmogorow-Axiomen lassen sich unmittelbar folgende Regeln ableiten:

\begin{satz}[Grundlegende Wahrscheinlichkeitsregeln]\label{satz:grundregeln}
  Sei $(\Omega, \mathcal{F}, \Prob)$ ein Wahrscheinlichkeitsraum und
  $A, B \in \mathcal{F}$. Dann gilt:
  \begin{enumerate}[label=(\roman*)]
    \item $\Prob(\emptyset) = 0$.
    \item $\Prob(A^c) = 1 - \Prob(A)$ \quad (Komplementregel).
    \item $A \subseteq B \Rightarrow \Prob(A) \leq \Prob(B)$ \quad (Monotonie).
    \item $\Prob(A \cup B) = \Prob(A) + \Prob(B) - \Prob(A \cap B)$ \quad
      (Siebformel / Inklusion-Exklusion).
    \item \textbf{Union Bound (Boole'sche Ungleichung):}
      $\Prob\!\left(\bigcup_{i=1}^{n} A_i\right) \leq \sum_{i=1}^{n} \Prob(A_i)$.
  \end{enumerate}
\end{satz}

\begin{proof}
  (i) Aus (K2) und (K3) folgt
  $1 = \Prob(\Omega) = \Prob(\Omega \cup \emptyset) = \Prob(\Omega) + \Prob(\emptyset)$,
  also $\Prob(\emptyset)=0$.\\
  (ii) $\Omega = A \cup A^c$ (disjunkt), daher $1 = \Prob(A) + \Prob(A^c)$.\\
  (iii) $B = A \cup (B \setminus A)$ (disjunkt), also $\Prob(B) = \Prob(A) + \Prob(B\setminus A) \geq \Prob(A)$.\\
  (iv) $A \cup B = A \cup (B \setminus A)$ (disjunkt), $\Prob(B\setminus A) = \Prob(B) - \Prob(A \cap B)$.
\end{proof}

\subsection{Laplace-Wahrscheinlichkeit}

Im klassischen Laplace-Modell ist $\Omega$ endlich und alle Elementarereignisse
sind gleich wahrscheinlich:

\begin{definition}[Laplace-Wahrscheinlichkeit]\label{def:laplace}
  Für ein endliches $\Omega$ mit $|\Omega| = n$ und
  $\Prob(\{\omega\}) = \frac{1}{n}$ für alle $\omega \in \Omega$ gilt für
  jedes Ereignis $A \subseteq \Omega$:
  \[
    \Prob(A) = \frac{|A|}{|\Omega|} = \frac{\text{Anzahl günstiger Ergebnisse}}
    {\text{Anzahl aller möglichen Ergebnisse}}.
  \]
\end{definition}

Die Vollständigkeitsbedingung ist hier trivialerweise erfüllt:
$\sum_{\omega\in\Omega} \frac{1}{|\Omega|} = |\Omega| \cdot \frac{1}{|\Omega|} = 1$.

\subsection{Unabhängigkeit von Ereignissen}

\begin{definition}[Stochastische Unabhängigkeit]\label{def:unabh}
  Zwei Ereignisse $A, B \in \mathcal{F}$ heißen \textbf{(stochastisch)
  unabhängig}, wenn
  \[
    \Prob(A \cap B) = \Prob(A) \cdot \Prob(B).
  \]
  Eine Familie $(A_i)_{i \in I}$ heißt \textbf{vollständig unabhängig}, wenn
  für jede endliche Teilmenge $J \subseteq I$:
  \[
    \Prob\!\left(\bigcap_{j \in J} A_j\right) = \prod_{j \in J} \Prob(A_j).
  \]
\end{definition}

% ===============================================================
\section{Bedingte und totale Wahrscheinlichkeit}
\label{sec:bedingt}
% ===============================================================

\subsection{Bedingte Wahrscheinlichkeit}

Ein fundamentales Konzept der Wahrscheinlichkeitstheorie ist die bedingte
Wahrscheinlichkeit: Wie verändert sich die Wahrscheinlichkeit eines Ereignisses,
wenn man über das Eintreten eines anderen Ereignisses informiert wird?

\begin{definition}[Bedingte Wahrscheinlichkeit]\label{def:bedingteW}
  Sei $(\Omega, \mathcal{F}, \Prob)$ ein Wahrscheinlichkeitsraum und
  $B \in \mathcal{F}$ mit $\Prob(B) > 0$. Die \textbf{bedingte
  Wahrscheinlichkeit} von $A$ gegeben $B$ ist definiert als
  \[
    \Prob(A \mid B) \coloneqq \frac{\Prob(A \cap B)}{\Prob(B)}.
  \]
\end{definition}

\begin{bemerkung}
  Die Abbildung $A \mapsto \Prob(A \mid B)$ ist selbst ein Wahrscheinlichkeitsmaß
  auf $(\Omega, \mathcal{F})$. Insbesondere gilt $\Prob(\Omega \mid B) = 1$, was
  die vollständige Beschreibbarkeit im bedingten Experiment sicherstellt.
\end{bemerkung}

\begin{beispiel}[Bedingte Wahrscheinlichkeit beim Würfel]\label{bsp:bedWuerfel}
  Beim einmaligen Würfelwurf sei $B = \{2, 4, 6\}$ (gerade Augenzahl) und
  $A = \{4, 5, 6\}$ (Augenzahl $\geq 4$). Dann:
  \[
    \Prob(A \mid B) = \frac{\Prob(A \cap B)}{\Prob(B)}
    = \frac{\Prob(\{4,6\})}{\Prob(\{2,4,6\})}
    = \frac{2/6}{3/6} = \frac{2}{3}.
  \]
  Wenn man weiß, dass eine gerade Zahl gefallen ist, steigt die Wahrscheinlichkeit
  für eine Zahl $\geq 4$ von $\frac{3}{6} = \frac{1}{2}$ auf $\frac{2}{3}$.
\end{beispiel}

\subsection{Multiplikationsregel und Kettenregel}

\begin{satz}[Multiplikationsregel]\label{satz:multregel}
  Für Ereignisse $A_1, \ldots, A_n \in \mathcal{F}$ mit
  $\Prob(A_1 \cap \cdots \cap A_{n-1}) > 0$ gilt:
  \[
    \Prob(A_1 \cap A_2 \cap \cdots \cap A_n)
    = \Prob(A_1) \cdot \Prob(A_2 \mid A_1)
      \cdot \Prob(A_3 \mid A_1 \cap A_2)
      \cdots \Prob(A_n \mid A_1 \cap \cdots \cap A_{n-1}).
  \]
\end{satz}

\subsection{Satz der totalen Wahrscheinlichkeit}

\begin{satz}[Satz der totalen Wahrscheinlichkeit]\label{satz:totalW}
  Sei $(B_i)_{i=1}^{n}$ eine Partition von $\Omega$ mit $\Prob(B_i) > 0$ für
  alle $i$. Dann gilt für jedes $A \in \mathcal{F}$:
  \[
    \Prob(A) = \sum_{i=1}^{n} \Prob(A \mid B_i) \cdot \Prob(B_i).
  \]
\end{satz}

\begin{proof}
  Da $(B_i)$ eine Partition ist, gilt $A = \bigsqcup_{i} (A \cap B_i)$.
  Mit $\sigma$-Additivität folgt
  $\Prob(A) = \sum_i \Prob(A \cap B_i) = \sum_i \Prob(A \mid B_i)\Prob(B_i)$.
\end{proof}

\subsection{Satz von Bayes}

Der Satz von Bayes ermöglicht die Umkehrung bedingter Wahrscheinlichkeiten
und ist das Fundament der bayesischen Statistik:

\begin{satz}[Satz von Bayes]\label{satz:bayes}
  Sei $(B_i)_{i=1}^{n}$ eine Partition von $\Omega$ mit $\Prob(B_i) > 0$ und
  $\Prob(A) > 0$. Dann gilt:
  \[
    \Prob(B_j \mid A)
    = \frac{\Prob(A \mid B_j) \cdot \Prob(B_j)}{\sum_{i=1}^{n} \Prob(A \mid B_i) \cdot \Prob(B_i)}.
  \]
  Man bezeichnet $\Prob(B_j)$ als \textbf{A-priori-Wahr\-schein\-lich\-keit} und
  $\Prob(B_j \mid A)$ als \textbf{A-pos\-te\-rio\-ri-Wahr\-schein\-lich\-keit}.
\end{satz}

Zentral ist, dass der Nenner gerade $\Prob(A) = 1$ sein würde, wenn $A = \Omega$,
was die Konsistenz mit dem Vollständigkeitsprinzip unterstreicht.

% ===============================================================
\section{Zufallsvariablen, Erwartungswert und Varianz}
\label{sec:erwartungswert}
% ===============================================================

\subsection{Zufallsvariablen}

Für Versuche, die durch diskrete Ereignisse beschrieben werden, ist es für
die Analyse manchmal hilfreich, eine Zufallsvariable zu verwenden. Dabei
beschreibt die Zufallsvariable den Zufall für das Ereignis variabel für den
konkreten Ausgang. Zum Beispiel kann für den Versuch, einen Würfel zu würfeln,
eine Zufallsvariable definiert werden durch
\begin{align}
  \label{math:var-x}
  X = \{\text{Anzahl der Augen, auf die der Würfel fällt}\}.
\end{align}
Es fiel auf, dass für jedes Ereignis dieses Versuches gilt, dass der Würfel
beim Wurf auf die Augenzahl $k$ fällt, die \textit{Wahrscheinlichkeit}
$\frac{1}{6}$ hat, $k \in \{1, \ldots, 6\}$. Mit der Zufallsvariable ist
diese Beobachtung nun so beschreibbar: $P(X = k) = \frac{1}{6}$, wobei $P$
für Probability steht.

\begin{definition}[Zufallsvariable]\label{def:zufvar}
  Eine \textbf{Zufallsvariable} (ZV) $X\colon \Omega \to \mathbb{R}$ ist eine
  messbare Funktion von einem Wahrscheinlichkeitsraum $(\Omega,\mathcal{F},\Prob)$
  nach $(\mathbb{R}, \mathcal{B}(\mathbb{R}))$. Die
  \textbf{Verteilungsfunktion} von $X$ ist
  \[
    F_X(x) \coloneqq \Prob(X \leq x), \quad x \in \mathbb{R}.
  \]
\end{definition}

Eine ZV $X$ heißt \textbf{diskret}, wenn sie nur abzählbar viele Werte annimmt;
sie heißt \textbf{stetig}, wenn $F_X$ absolut stetig ist, d.\,h.\ eine
Dichte $f_X$ besitzt mit $F_X(x) = \int_{-\infty}^{x} f_X(t)\,dt$.

\subsection{Erwartungswert}

Bei dem Versuch, den Würfel zu würfeln, ist es vor dem Versuch von Bedeutung,
zu ermitteln, welcher Ausgang des Versuches zu erwarten ist. Für Versuche,
die durch diskrete Ereignisse beschrieben werden, wird der Erwartungswert
für eine Zufallsvariable $X$ berechnet.

\begin{definition}[Erwartungswert]\label{def:ew}
  Sei $X$ eine ZV auf $(\Omega, \mathcal{F}, \Prob)$.
  \begin{itemize}
    \item \textbf{Diskreter Fall:} $X$ nimmt die Werte $x_1, x_2, \ldots$ an.
      Existiert $\sum_i |x_i| \Prob(X = x_i) < \infty$, so ist
      \[
        \E[X] \coloneqq \sum_{i} x_i \cdot \Prob(X = x_i).
      \]
    \item \textbf{Stetiger Fall:} $X$ besitzt Dichte $f_X$. Existiert
      $\int_{-\infty}^{\infty} |x| f_X(x)\,dx < \infty$, so ist
      \[
        \E[X] \coloneqq \int_{-\infty}^{\infty} x \cdot f_X(x)\,dx.
      \]
    \item \textbf{Allgemeiner Fall (Lebesgue-Integral):}
      $\E[X] \coloneqq \int_\Omega X(\omega)\,d\Prob(\omega)$.
  \end{itemize}
\end{definition}

\begin{beispiel}[Erwartungswert beim fairen Würfel]\label{bsp:ewwuerfel}
  $X$ sei die Augenzahl beim einmaligen Würfeln gemäß \eqref{math:var-x},
  $\Prob(X=k)=\frac{1}{6}$ für $k \in \{1,\ldots,6\}$. Dann:
  \begin{align}
    E[X] &= \sum_{k \in \{1, \ldots, 6\}} k \cdot P(X = k) \\
         &= 1 \cdot P(X = 1) + \ldots + 6 \cdot P(X = 6) \\
         &= 1 \cdot \frac{1}{6} + \ldots + 6 \cdot \frac{1}{6} \\
         &= 3{,}5.
  \end{align}
  Das bedeutet, bevor wir den Würfel einmal gewürfelt haben, erwarten wir,
  dass der Würfel auf die Augenzahl $3{,}5$ fällt. Diese Augenzahl gibt es
  nicht. Der Erwartungswert aber beschreibt formal am genauesten, welcher
  Ausgang bei diesem Versuch wirklich zu erwarten ist. Denn erst wenn wir
  diesen Versuch oft wiederholen, ist es interessant, wie der tatsächliche
  Ausgang des Versuches vom Erwartungswert abweicht. Beachte:
  $3{,}5 \notin \Omega$. Der Erwartungswert ist der \emph{langfristige
  Durchschnitt}, keine einzelne Augenzahl.
\end{beispiel}

\begin{satz}[Linearität des Erwartungswerts]\label{satz:linear}
  Für ZV $X, Y$ und Konstanten $a, b \in \mathbb{R}$ gilt stets:
  \[
    \E[aX + bY] = a\,\E[X] + b\,\E[Y].
  \]
  Diese Eigenschaft gilt \textbf{unabhängig} davon, ob $X$ und $Y$ unabhängig
  sind.
\end{satz}

Die Linearität des Erwartungswerts ist eines der mächtigsten Werkzeuge der
Wahrscheinlichkeitstheorie und fundamental für die Analyse randomisierter
Algorithmen.

\subsection{Varianz und Standardabweichung}

\begin{definition}[Varianz und Standardabweichung]\label{def:var}
  \[
    \Var(X) \coloneqq \E\!\left[(X - \E[X])^2\right]
    = \E[X^2] - (\E[X])^2.
  \]
  Die \textbf{Standardabweichung} ist $\sigma_X \coloneqq \sqrt{\Var(X)}$.
\end{definition}

\begin{satz}[Rechenregeln für die Varianz]\label{satz:varregeln}
  \begin{enumerate}[label=(\roman*)]
    \item $\Var(aX + b) = a^2 \Var(X)$ für $a, b \in \mathbb{R}$.
    \item Sind $X_1, \ldots, X_n$ paarweise unkorreliert (insb.\ unabhängig):
      \[
        \Var\!\left(\sum_{i=1}^n X_i\right) = \sum_{i=1}^n \Var(X_i).
      \]
  \end{enumerate}
\end{satz}

\subsection{Wichtige Ungleichungen}

\begin{satz}[Markov-Ungleichung]\label{satz:markov}
  Für eine nicht-negative ZV $X$ und $a > 0$:
  \[
    \Prob(X \geq a) \leq \frac{\E[X]}{a}.
  \]
\end{satz}

\begin{satz}[Tschebyschow-Ungleichung]\label{satz:tscheb}
  Für eine ZV $X$ mit $\E[X] = \mu$ und $\Var(X) = \sigma^2$ sowie $k > 0$:
  \[
    \Prob\!\left(|X - \mu| \geq k\sigma\right) \leq \frac{1}{k^2}.
  \]
  Äquivalent: $\Prob(|X-\mu| \geq t) \leq \frac{\sigma^2}{t^2}$.
\end{satz}

\begin{satz}[Chernoff-Schranken]\label{satz:chernoff}
  Seien $X_1, \ldots, X_n$ unabhängige Bernoulli-ZV mit $\Prob(X_i=1) = p_i$
  und $\mu = \E\!\left[\sum X_i\right] = \sum p_i$. Dann gilt für $\delta > 0$:
  \[
    \Prob\!\left(\sum_{i=1}^n X_i \geq (1+\delta)\mu\right)
    \leq \left(\frac{e^\delta}{(1+\delta)^{1+\delta}}\right)^{\mu}
    \leq e^{-\mu\delta^2/3} \quad \text{für } \delta \leq 1.
  \]
  Für die untere Schranke gilt für $0 < \delta < 1$:
  \[
    \Prob\!\left(\sum_{i=1}^n X_i \leq (1-\delta)\mu\right)
    \leq e^{-\mu\delta^2/2}.
  \]
\end{satz}

\subsection{Bedingter Erwartungswert}

\begin{definition}[Bedingter Erwartungswert]\label{def:bedEW}
  Sei $X$ eine ZV und $B \in \mathcal{F}$ mit $\Prob(B) > 0$. Der
  \textbf{bedingte Erwartungswert} von $X$ gegeben $B$ ist
  \[
    \E[X \mid B] \coloneqq \frac{\E[X \cdot \mathbf{1}_B]}{\Prob(B)}
    = \frac{1}{\Prob(B)} \int_B X(\omega)\,d\Prob(\omega).
  \]
  Für eine diskrete ZV $Y$ mit Werten $y_1, y_2, \ldots$ ist der
  \textbf{bedingte Erwartungswert} von $X$ gegeben $Y$ die ZV
  $\E[X \mid Y] \coloneqq g(Y)$, wobei
  $g(y_j) = \E[X \mid Y = y_j]$.
\end{definition}

\begin{satz}[Turm-Eigenschaft (Law of Total Expectation)]\label{satz:turm}
  \[
    \E\!\left[\E[X \mid Y]\right] = \E[X].
  \]
  Insbesondere: Ist $(B_i)$ eine Partition von $\Omega$ mit $\Prob(B_i)>0$,
  so gilt
  \[
    \E[X] = \sum_i \E[X \mid B_i] \cdot \Prob(B_i).
  \]
\end{satz}

Die Turm-Eigenschaft ist das Analogon des Satzes der totalen
Wahrscheinlichkeit für Erwartungswerte. Sie drückt erneut das Prinzip der
vollständigen Beschreibbarkeit aus: Man kann den Erwartungswert berechnen,
indem man über eine vollständige Partition des Ergebnisraums mittelt.

% ===============================================================
\section{Diskrete Wahrscheinlichkeitsverteilungen}
\label{sec:verteilungen}
% ===============================================================

\subsection{Überblick und Vollständigkeitsprüfung}

Wann immer die Lehre vom Zufall in der Mathematik so verlassen wird, dass sie
in der Praxis nicht mehr anwendbar ist, verletzen wir das eigentliche Ziel der
Mathematik in diesem wichtigen Thema. Beim Erwartungswert war dies schon der
Fall, denn auf die Augenzahl $3{,}5$ wird der Würfel nie zufällig fallen.
Damit muss die Lehre vom Zufall in der Mathematik sehr vorsichtig entwickelt
werden. Denn wenn man zum Beispiel den Versuch, einen Würfel zu würfeln, oft
wiederholt, ist beobachtbar, dass die Abweichung vom Erwartungswert beachtlich
ist.

Die bekannten Wahrscheinlichkeitsverteilungen der Mathematik sind für die
unterschiedlichen Versuche mit diskreten Ereignissen mit großer Genauigkeit
und Vorsicht anzuwenden. Wird ein Versuch in seiner Beschreibung auf eine
Anwendung einer Wahrscheinlichkeitsverteilung überprüft, können kleinste
Details in der Beschreibung des Versuchs entscheidend sein für die Wahl der
Wahrscheinlichkeitsverteilung. Wird ein Detail in der Beschreibung des
Versuchs übersehen und eine andere Wahrscheinlichkeitsverteilung gewählt,
kann dies in der Analyse für den Versuch zu einem anderen Ergebnis führen.

Im Folgenden werden die unterschiedlichen Wahrscheinlichkeitsverteilungen
aufgeführt. Für jede Verteilung wird explizit gezeigt, dass die
Vollständigkeitsbedingung $\sum_k \Prob(X=k) = 1$ erfüllt ist. Dabei sollen
Gemeinsamkeiten auffallen und dazu führen, dass diese Verteilungen mit
Genauigkeit und Vorsicht anzuwenden sind.

\subsection{Bernoulli-Verteilung}

Die Bernoulli-Verteilung beschreibt Versuche, bei denen es nur zwei mögliche
Ereignisse gibt. Ein Ereignis wird als Erfolg bezeichnet und das andere als
Misserfolg. Zum Beispiel kann beim Wurf einer Münze das Ereignis Kopf als
Erfolg und das Ereignis Zahl als Misserfolg bezeichnet werden. Ebenso kann
beim Würfeln eines Würfels das Ereignis, dass der Würfel auf die Augenzahl
$6$ fällt, als Erfolg bezeichnet werden und alle anderen Ereignisse, dass der
Würfel auf die Augenzahlen $1, 2, 3, 4$ oder $5$ fällt, als Misserfolg.

\begin{definition}[Bernoulli-Verteilung $\operatorname{Ber}(p)$]\label{def:bernoulli}
  Für einen Versuch, der durch die Bernoulli-Verteilung beschrieben wird,
  wird eine Zufallsvariable $X$ definiert durch
  \begin{align}
    X =
    \begin{cases}
      1, & \text{falls Erfolg eintritt} \\
      0, & \text{falls Misserfolg eintritt}
    \end{cases}
  \end{align}
  Die Wahrscheinlichkeit für den Erfolg wird mit $p$ bezeichnet, also
  $P(X = 1) = p$. Die Wahrscheinlichkeit für den Misserfolg ist dann
  $P(X = 0) = 1 - p$. Es ist wichtig zu beachten, dass die Summe aller
  Wahrscheinlichkeiten $p + (1-p) = 1$ ergibt. Damit ist der Versuch
  mathematisch vollständig beschrieben. Es gilt
  $\E[X] = p$ und $\Var(X) = p(1-p)$.
\end{definition}

Der Erwartungswert einer bernoulli-verteilten Zufallsvariablen $X$ berechnet
sich wie folgt:
\begin{align}
  E[X] &= 1 \cdot P(X = 1) + 0 \cdot P(X = 0) \\
       &= 1 \cdot p + 0 \cdot (1-p) \\
       &= p.
\end{align}
Das bedeutet, bei einem Versuch, der durch die Bernoulli-Verteilung beschrieben
wird, erwarten wir im Mittel den Wert $p$. Für den Münzwurf mit
$p = \frac{1}{2}$ erwarten wir also den Wert $0{,}5$. Dieser Wert kann nicht
eintreten, denn die Münze zeigt entweder Kopf oder Zahl. Der Erwartungswert
beschreibt aber formal genau, was bei vielen Wiederholungen dieses Versuchs
zu erwarten ist.

\begin{beispiel}[Basketball]\label{bsp:basketball}
  Ein Basketballspieler trifft den Korb mit einer Wahrscheinlichkeit von
  $p = 0{,}7$. Bei einem einzelnen Wurf wird der Erfolg (Treffer) mit der
  Zufallsvariablen $X$ beschrieben, die bernoulli-verteilt ist mit Parameter
  $p = 0{,}7$. Der erwartete Wert ist $E[X] = 0{,}7$.
\end{beispiel}

Die Bernoulli-Verteilung ist die einfachste aller Wahrscheinlichkeitsverteilungen.
Sie ist die Grundlage für viele weitere Verteilungen. Wird ein Versuch, der
durch die Bernoulli-Verteilung beschrieben wird, mehrfach unabhängig
wiederholt, ergeben sich andere Wahrscheinlichkeitsverteilungen.

\subsection{Binomialverteilung}

\begin{definition}[Binomialverteilung $\operatorname{Bin}(n,p)$]\label{def:binom}
  $X \sim \operatorname{Bin}(n,p)$: Anzahl der Erfolge bei $n$ unabhängigen
  Bernoulli-Experimen\-ten mit Erfolgswahrscheinlichkeit $p$.
  \[
    \Prob(X=k) = \binom{n}{k} p^k (1-p)^{n-k}, \quad k=0,\ldots,n.
  \]
  Vollständigkeit (Binomischer Lehrsatz):
  $\sum_{k=0}^{n}\binom{n}{k}p^k(1-p)^{n-k} = (p+(1-p))^n = 1$. \\
  $\E[X] = np$, $\Var(X) = np(1-p)$.
\end{definition}

\subsection{Geometrische Verteilung}

Die geometrische Verteilung beschreibt Versuche, bei denen ein
Bernoulli-Versuch so oft wiederholt wird, bis zum ersten Mal ein Erfolg
eintritt. Es stellt sich die Frage: Wie viele Versuche sind nötig, bis der
erste Erfolg eintritt? Beim Würfeln eines Würfels kann man fragen: Wie oft
muss der Würfel geworfen werden, bis zum ersten Mal die Augenzahl $6$ fällt?
Bei diesem Versuch wird der Würfel so lange geworfen, bis die Augenzahl $6$
das erste Mal erscheint.

\begin{definition}[Geometrische Verteilung $\operatorname{Geo}(p)$]\label{def:geo}
  Für einen Versuch, der durch die geometrische Verteilung beschrieben wird,
  wird eine Zufallsvariable $X$ definiert durch
  \begin{align}
    X = \{\text{Anzahl der Versuche bis zum ersten Erfolg}\}.
  \end{align}
  Die Wahrscheinlichkeit, dass der erste Erfolg genau beim $k$-ten Versuch
  eintritt: Es müssen zunächst $k-1$ Misserfolge eintreten, jeder mit
  Wahrscheinlichkeit $1-p$, und dann muss beim $k$-ten Versuch ein Erfolg
  mit Wahrscheinlichkeit $p$ eintreten. Damit ergibt sich:
  \[
    \Prob(X=k) = (1-p)^{k-1} p, \quad k = 1, 2, \ldots
  \]
  Die Vollständigkeit überprüfen wir explizit:
  \begin{align}
    \sum_{k=1}^{\infty} P(X = k)
    = \sum_{k=1}^{\infty} (1-p)^{k-1} \cdot p
    = p \cdot \frac{1}{1-(1-p)} = p \cdot \frac{1}{p} = 1.
  \end{align}
  Damit ist der Versuch mathematisch vollständig beschrieben.
  $\E[X] = \frac{1}{p}$, $\Var(X) = \frac{1-p}{p^2}$.
\end{definition}

Das bedeutet, im Erwartungswert benötigen wir $\frac{1}{p}$ Versuche bis zum
ersten Erfolg. Für das Würfeln einer $6$ mit einem Würfel ist
$p = \frac{1}{6}$. Damit erwarten wir $E[X] = \frac{1}{1/6} = 6$ Würfe bis
zum ersten Mal die Augenzahl $6$ fällt. Dies ist ein intuitives Ergebnis.
Wird der Versuch in der Praxis durchgeführt, kann es sein, dass die $6$
bereits beim ersten Wurf fällt. Es kann aber auch sein, dass die $6$ erst
nach 20 Würfen das erste Mal fällt. Der Erwartungswert beschreibt das, was
im Mittel über viele Wiederholungen dieses gesamten Versuchs zu erwarten ist.

\begin{beispiel}[Telefonverkauf]\label{bsp:telefonverkauf}
  Ein Unternehmen ruft potenzielle Kunden an. Jeder Kunde kauft mit
  Wahrscheinlichkeit $p = 0{,}2$ das Produkt. Die Anzahl der Anrufe bis zum
  ersten Verkauf ist geometrisch verteilt mit Parameter $p = 0{,}2$. Im
  Erwartungswert sind $E[X] = \frac{1}{0{,}2} = 5$ Anrufe nötig bis zum
  ersten Verkauf.
\end{beispiel}

Die geometrische Verteilung zeigt, wie vorsichtig der Erwartungswert
interpretiert werden muss. In der Praxis kann die tatsächliche Anzahl der
benötigten Versuche erheblich vom Erwartungswert abweichen. Dies ist bei allen
Wahrscheinlichkeitsverteilungen zu beachten.

\subsection{Poisson-Verteilung}

\begin{definition}[Poisson-Verteilung $\operatorname{Poi}(\lambda)$]\label{def:poisson}
  \[
    \Prob(X=k) = \frac{\lambda^k e^{-\lambda}}{k!}, \quad k=0,1,2,\ldots,\;
    \lambda > 0.
  \]
  Vollständigkeit (Exponentialreihe):
  $\sum_{k=0}^{\infty}\frac{\lambda^k e^{-\lambda}}{k!}
   = e^{-\lambda} \sum_{k=0}^{\infty}\frac{\lambda^k}{k!}
   = e^{-\lambda} \cdot e^{\lambda} = 1$. \\
  $\E[X] = \lambda$, $\Var(X) = \lambda$.
\end{definition}

\subsection{Stetige Verteilungen}

\subsubsection{Gleichverteilung $\mathcal{U}([a,b])$}

\[
  f(x) = \frac{1}{b-a} \cdot \mathbf{1}_{[a,b]}(x).
\]
Vollständigkeit: $\int_a^b \frac{1}{b-a}\,dx = 1$.
$\E[X] = \frac{a+b}{2}$, $\Var(X) = \frac{(b-a)^2}{12}$.

\subsubsection{Normalverteilung $\mathcal{N}(\mu, \sigma^2)$}

Die Gaußsche Verteilung, auch Normalverteilung genannt, beschreibt
kontinuierliche Ereignisse. Bei kontinuierlichen Ereignissen kann die
Zufallsvariable jeden beliebigen Wert aus einem Intervall annehmen, zum
Beispiel alle reellen Zahlen. Ein Beispiel für einen Versuch mit
kontinuierlichen Ereignissen ist die Messung der Körpergröße von Menschen.
Die Körpergröße kann jeden Wert aus einem bestimmten Intervall annehmen –
sie ist keine diskrete Größe wie die Augenzahl beim Würfeln, sondern eine
kontinuierliche Größe.

Die Gaußsche Verteilung wird durch zwei Parameter beschrieben: den
Erwartungswert $\mu$ und die Varianz $\sigma^2$. Der Erwartungswert $\mu$
gibt an, um welchen Wert herum die Zufallsvariable streut. Die Varianz
$\sigma^2$ gibt an, wie stark die Zufallsvariable um den Erwartungswert
streut. Die Standardabweichung $\sigma$ ist die Wurzel aus der Varianz.

\[
  f(x) = \frac{1}{\sigma\sqrt{2\pi}}
  \exp\!\left(-\frac{(x-\mu)^2}{2\sigma^2}\right), \quad x \in \mathbb{R}.
\]
Bei kontinuierlichen Verteilungen kann nicht mehr nach der Wahrscheinlichkeit
für ein einzelnes Ereignis gefragt werden, denn diese ist immer $0$. Stattdessen
wird nach der Wahrscheinlichkeit gefragt, dass die Zufallsvariable $X$ in
einem bestimmten Intervall $[a, b]$ liegt:
\[
  \Prob(a \le X \le b) = \int_a^b f(x)\,dx.
\]
Vollständigkeit: $\int_{-\infty}^{\infty} f(x)\,dx = 1$ (Gaußsches Integral).
$\E[X] = \mu$, $\Var(X) = \sigma^2$.

\begin{beispiel}[Körpergröße]\label{bsp:koerpergroesse}
  Die Körpergröße von erwachsenen Männern in Deutschland ist annähernd
  normalverteilt mit Erwartungswert $\mu = 180$\,cm und Standardabweichung
  $\sigma = 7$\,cm. Die Wahrscheinlichkeit, dass ein zufällig ausgewählter
  Mann eine Körpergröße zwischen $173$\,cm und $187$\,cm hat, entspricht
  etwa $68\,\%$. Dies ist ein bekanntes Ergebnis der Gaußschen Verteilung:
  Etwa $68\,\%$ aller Werte liegen im Intervall $[\mu - \sigma, \mu +
  \sigma]$.
\end{beispiel}

Nicht jede Größe ist normalverteilt. Wird eine Größe fälschlicherweise als
normalverteilt angenommen, können die Ergebnisse der Analyse völlig falsch
sein. Zum Beispiel sind Aktienkurse nicht normalverteilt, obwohl dies oft
angenommen wird. Extreme Ereignisse treten bei Aktienkursen viel häufiger auf
als die Gaußsche Verteilung vorhersagt. Die Gaußsche Verteilung ist ein
mächtiges Werkzeug, aber sie muss mit Genauigkeit und Vorsicht angewendet
werden. Nur wenn die Voraussetzungen für ihre Anwendung erfüllt sind, liefert
sie zuverlässige Ergebnisse.

\subsubsection{Exponentialverteilung $\operatorname{Exp}(\lambda)$}

\[
  f(x) = \lambda e^{-\lambda x} \cdot \mathbf{1}_{[0,\infty)}(x), \quad \lambda>0.
\]
Vollständigkeit: $\int_0^{\infty}\lambda e^{-\lambda x}\,dx = 1$.
$\E[X] = \frac{1}{\lambda}$, $\Var(X) = \frac{1}{\lambda^2}$.

Die Exponentialverteilung besitzt die \textbf{Gedächtnislosigkeit}:
$\Prob(X > s+t \mid X > s) = \Prob(X > t)$.

\subsection{Anwendung: Hybrider Registerspeicher}
\label{subsec:registerspeicher}

Ein interessantes Beispiel der Gaußschen Verteilung ist die Optimierung von
Registerspeichern in Prozessoren. Dabei wird beobachtet, dass die Zugriffe
auf Register nicht gleichmäßig verteilt sind, sondern einer Gaußschen
Verteilung folgen. Einige Register werden sehr häufig verwendet, während
andere Register selten genutzt werden.

Anstatt alle Register gleich zu behandeln, wird eine hierarchische Struktur
eingeführt. Die Register, die am häufigsten genutzt werden, werden in
schnellem und teurem Speicher realisiert. Die Register, die seltener genutzt
werden, werden in langsamem und günstigem Speicher realisiert.

Für diesen Versuch wird eine Zufallsvariable $X$ definiert durch
\[
  X = \{\text{Register-Index, auf den zugegriffen wird}\}.
\]
Die Beobachtung zeigt, dass $X$ annähernd normalverteilt ist mit einem
Erwartungswert $\mu$ und einer Standardabweichung $\sigma$. Der Erwartungswert
$\mu$ gibt an, welcher Register-Index im Mittel am häufigsten genutzt wird.

Der hierarchische Aufbau des Registerspeichers erfolgt in zwei Bereichen. Der
erste Bereich, genannt \emph{Tier~1}, enthält etwa 32 bis 64 Register, die
nahe beim Erwartungswert $\mu$ liegen, also im Intervall
$[\mu - \sigma, \mu + \sigma]$. Wie beim Körpergröße-Beispiel wissen wir,
dass etwa $68\,\%$ aller Zugriffe auf Register in diesem Intervall liegen.
Diese Register werden in schnellem Speicher realisiert. Der zweite Bereich,
genannt \emph{Tier~2}, enthält alle übrigen Register in langsamem, aber
günstigem Speicher.

Um zu ermitteln, welche Register zu Tier~1 gehören, muss das Zugriffsmuster
analysiert werden. Für jedes Register $i$ wird während des Betriebs eine
Zufallsvariable $Y_i$ gezählt (Anzahl der Zugriffe auf Register~$i$). Die
Register mit den höchsten Werten von $Y_i$ werden Tier~1 zugeordnet.

Es ist wichtig zu beachten, dass die Zuordnung nicht statisch ist – die
Zugriffsmuster können sich während der Laufzeit ändern. Tier~1 und Tier~2
müssen daher dynamisch angepasst werden.

Die Kostenersparnis durch diesen Ansatz lässt sich berechnen. Sei $n_1$ die
Anzahl der Register in Tier~1 und $n_2$ die Anzahl in Tier~2, $c_1 \gg c_2$
die jeweiligen Kosten pro Register. Die Gesamtkosten betragen
$C = n_1 c_1 + n_2 c_2$. Ohne die hierarchische Struktur wären alle
$n_1 + n_2$ Register in teurem Speicher mit Gesamtkosten
$(n_1 + n_2) c_1$. Die Kostenersparnis beträgt damit
\begin{align}
  \Delta C = (n_1 + n_2)\,c_1 - (n_1\,c_1 + n_2\,c_2) = n_2\,(c_1 - c_2).
\end{align}
Da typischerweise $n_2 \gg n_1$ gilt, ist die Ersparnis erheblich. Dieses
Beispiel zeigt, wie die Gaußsche Verteilung in der Praxis angewendet wird –
aber auch, wie wichtig es ist, die Voraussetzung der Normalverteilung für
das konkrete Programm zu überprüfen.

% ===============================================================
\section{Randomisierte Algorithmen}
\label{sec:algorithmen}
% ===============================================================

Zufall ist nicht nur ein Gegenstand mathematischer Analyse, sondern ein
mächtiges \emph{algorithmisches Werkzeug}. Randomisierte Algorithmen nutzen
Zufallsentscheidungen, um
\begin{itemize}
  \item im \textbf{Worst Case schwere} Probleme \emph{im Erwartungswert}
    effizient zu lösen (\emph{Las-Vegas-Algorithmen}: immer korrekt,
    erwartete Laufzeit gut),
  \item korrekte Lösungen \emph{mit hoher Wahrscheinlichkeit} schnell zu
    finden (\emph{Monte-Carlo-Algorithmen}: eventuell falsch, aber mit
    kleiner Fehlerwahrscheinlichkeit),
  \item im kontinuierlichen Raum integrierbar zu machen
    (\emph{Monte-Carlo-Integration}).
\end{itemize}

\subsection{3-SAT-Approximation mit Erwartungswert $\geq 7/8$}
\label{subsec:3sat}

\subsubsection{Das Problem}

\begin{definition}[3-SAT]\label{def:3sat}
  Gegeben eine aussagenlogische Formel $\varphi = C_1 \wedge C_2 \wedge \cdots
  \wedge C_m$ in konjunktiver Normalform, wobei jede Klausel $C_j$ genau drei
  Literale enthält. Gesucht ist eine Variablenbelegung, die so viele Klauseln
  wie möglich erfüllt. Das Entscheidungsproblem (alle Klauseln erfüllbar?) ist
  NP-vollständig \cite{cook1971complexity}.
\end{definition}

\subsubsection{Der randomisierte Algorithmus}

Die folgende Idee stammt von Johnson \cite{johnson1974approximation}:
Weise jeder Variable $x_i$ unabhängig und gleichverteilt einen der Werte
$\{\mathtt{wahr}, \mathtt{falsch}\}$ zu.

\begin{algorithm}[H]
  \caption{Randomisierte 3-SAT-Approximation}
  \label{alg:3sat}
  \begin{algorithmic}[1]
    \Require{3-CNF-Formel $\varphi$ mit Variablen $x_1, \ldots, x_n$ und
             Klauseln $C_1, \ldots, C_m$}
    \Ensure{Variablenbelegung $\beta$ (erfüllt in Erwartung $\geq 7m/8$ Klauseln)}
    \For{$i \gets 1$ \textbf{to} $n$}
      \State Wähle $\beta(x_i) \in \{\texttt{wahr}, \texttt{falsch}\}$
             gleichverteilt und unabhängig.
    \EndFor
    \State $\mathit{SAT} \gets 0$
    \For{$j \gets 1$ \textbf{to} $m$}
      \If{$C_j$ ist unter $\beta$ erfüllt}
        \State $\mathit{SAT} \gets \mathit{SAT} + 1$
      \EndIf
    \EndFor
    \State \Return $\beta$ (und $\mathit{SAT}$)
  \end{algorithmic}
\end{algorithm}

\subsubsection{Erwartungswert-Analyse}

\begin{satz}[Erwartungswert der erfüllten Klauseln]\label{satz:ewsat}
  Sei $Y_j = \mathbf{1}[\text{Klausel } C_j \text{ ist erfüllt}]$ für
  $j = 1, \ldots, m$. Dann gilt:
  \[
    \E\!\left[\sum_{j=1}^m Y_j\right] \geq \frac{7}{8}\, m.
  \]
\end{satz}

\begin{proof}
  Da jede Klausel genau drei Literale enthält, ist eine Klausel genau dann
  \emph{nicht} erfüllt, wenn alle drei Literale den Wert $\mathtt{falsch}$
  haben. Jedes Literal ist unabhängig mit Wahrscheinlichkeit $\frac{1}{2}$
  falsch. Also:
  \[
    \Prob(C_j \text{ nicht erfüllt}) = \left(\frac{1}{2}\right)^{\!3} = \frac{1}{8}.
  \]
  Daher:
  \[
    \Prob(C_j \text{ erfüllt}) = 1 - \frac{1}{8} = \frac{7}{8}.
  \]
  Mit der Linearität des Erwartungswerts (Satz~\ref{satz:linear}):
  \[
    \E\!\left[\sum_{j=1}^m Y_j\right] = \sum_{j=1}^m \E[Y_j]
    = \sum_{j=1}^m \Prob(C_j \text{ erfüllt}) = m \cdot \frac{7}{8}.
  \]
  Hierbei wurde \emph{keine} Unabhängigkeit der $Y_j$ benötigt – nur die
  Linearität des Erwartungswerts.
\end{proof}

Die erwartete Gesamtzahl aller erfüllten Klauseln beträgt damit also
$\frac{7}{8}m$. Sei $O$ die optimale Anzahl an Klauseln, die in $\varphi$
erfüllt werden können. Dann kann $O$ höchstens alle $m$ Klauseln erfüllen,
also $O \le m$. Damit ergibt sich
\[
  \E[Y] = \frac{7}{8}m \ge \frac{7}{8}\, O.
\]
Das bedeutet: Die Lösung des Algorithmus erfüllt im Erwartungswert mindestens
$\frac{7}{8}$ aller Klauseln der optimalen Lösung, die überhaupt in $\varphi$
erfüllt werden können. Die Laufzeit ist linear in der Anzahl der Klauseln
und Variablen.

\paragraph{Verbindung zum Vollständigkeitsprinzip.}
Der Beweis nutzt implizit die vollständige Beschreibbarkeit: Für jede Klausel
$C_j$ ist $\Prob(C_j \text{ erfüllt}) + \Prob(C_j \text{ nicht erfüllt})
= \frac{7}{8}+\frac{1}{8}=1$.

\subsubsection{Derandomisierung via Methode der bedingten Erwartungswerte}

Mit der Methode der bedingten Erwartungswerte \cite{alon2016probabilistic}
lässt sich der Algorithmus \emph{derandomisieren}: Man setzt die Variablen
nacheinander auf $\mathtt{wahr}$ oder $\mathtt{falsch}$, indem man bei
jedem Schritt denjenigen Wert wählt, der den bedingten Erwartungswert der
erfüllten Klauseln maximiert. Das resultierende deterministische Verfahren
hat Laufzeit $O(nm)$ und garantiert ebenfalls $\geq \frac{7}{8}m$ erfüllte
Klauseln. HÃ¥stad zeigte, dass $\frac{7}{8}$ optimal unter der Annahme
$\text{P} \neq \text{NP}$ ist \cite{hastad2001some}.

\subsubsection{Laufzeitanalyse}

\begin{itemize}
  \item Schleife zur Belegung: $\Theta(n)$ Zufallszüge.
  \item Auswertung: $\Theta(m)$ (jede Klausel hat konstante Größe~3).
  \item Gesamtlaufzeit: $\Theta(n + m)$ – linear in der Eingabegröße.
\end{itemize}

\subsection{Randomisiertes Quicksort}
\label{subsec:rqs}

\subsubsection{Idee und Pseudocode}

Beim deterministischen Quicksort kann ein adversarial gewähltes Pivot-Element
zur Worst-Case-Laufzeit $\Theta(n^2)$ führen. Die randomisierte Variante
wählt das Pivot-Element zufällig und gleichverteilt:

\begin{algorithm}[H]
  \caption{Randomisiertes Quicksort}
  \label{alg:rqs}
  \begin{algorithmic}[1]
    \Require{Array $A[1..n]$ mit $n$ verschiedenen Elementen}
    \Ensure{$A$ sortiert in aufsteigender Reihenfolge}
    \Function{RandQuicksort}{$A, l, r$}
      \If{$l \geq r$}
        \State \Return
      \EndIf
      \State $p \gets \textsc{Zufälliger-Index}(l, r)$
        \Comment{gleichverteilt in $\{l, \ldots, r\}$}
      \State Tausche $A[p]$ mit $A[r]$
      \State $q \gets \textsc{Partition}(A, l, r)$
        \Comment{Lomuto- oder Hoare-Partition}
      \State \Call{RandQuicksort}{$A, l, q-1$}
      \State \Call{RandQuicksort}{$A, q+1, r$}
    \EndFunction
    \Statex
    \Function{Partition}{$A, l, r$}
      \State $\mathit{pivot} \gets A[r]$
      \State $i \gets l - 1$
      \For{$j \gets l$ \textbf{to} $r-1$}
        \If{$A[j] \leq \mathit{pivot}$}
          \State $i \gets i+1$
          \State Tausche $A[i]$ mit $A[j]$
        \EndIf
      \EndFor
      \State Tausche $A[i+1]$ mit $A[r]$
      \State \Return $i+1$
    \EndFunction
  \end{algorithmic}
\end{algorithm}

\subsubsection{Erwartete Laufzeitanalyse}

\begin{satz}[Erwartete Laufzeit von Randomisiertem Quicksort]\label{satz:ewrqs}
  Die erwartete Anzahl von Vergleichen beim randomisierten Quicksort
  auf einem Array der Länge $n$ (mit paarweise verschiedenen Elementen)
  beträgt:
  \[
    \E[T(n)] = 2n\ln n + O(n) = O(n \log n).
  \]
\end{satz}

\begin{proof}
  Seien $z_1 < z_2 < \cdots < z_n$ die sortierten Elemente. Definiere die
  Indikatorvariable
  \[
    X_{ij} = \mathbf{1}[\text{$z_i$ und $z_j$ werden verglichen}],
    \quad 1 \leq i < j \leq n.
  \]
  Die Gesamtzahl der Vergleiche ist $C = \sum_{i<j} X_{ij}$.
  Mit der Linearität des Erwartungswerts:
  \[
    \E[C] = \sum_{i=1}^{n}\sum_{j=i+1}^{n} \Prob(z_i \text{ und } z_j \text{ werden verglichen}).
  \]
  \textbf{Schlüsselbeobachtung:} $z_i$ und $z_j$ werden genau dann verglichen,
  wenn in der Menge $\{z_i, z_{i+1}, \ldots, z_j\}$ entweder $z_i$ oder
  $z_j$ als erstes zum Pivot gewählt wird. Da das Pivot gleichverteilt aus
  dieser Menge der $j-i+1$ Elemente gewählt wird:
  \[
    \Prob(X_{ij} = 1) = \frac{2}{j - i + 1}.
  \]
  Vollständigkeit: Es gilt
  \begin{align*}
    &\Prob(\text{$z_i$ zuerst Pivot in $\{z_i,\ldots,z_j\}$})
    + \Prob(\text{$z_j$ zuerst Pivot})
    + \Prob(\text{anderes Element zuerst}) \\
    &\quad= \frac{1}{j{-}i{+}1} + \frac{1}{j{-}i{+}1} + \frac{j{-}i{-}1}{j{-}i{+}1} = 1.
  \end{align*}

  Damit ergibt sich:
  \begin{align*}
    \E[C] &= \sum_{i=1}^{n}\sum_{j=i+1}^{n} \frac{2}{j-i+1}
    = \sum_{i=1}^{n}\sum_{k=2}^{n-i+1} \frac{2}{k}
    \leq \sum_{i=1}^{n}\sum_{k=2}^{n} \frac{2}{k} \\
    &= 2n \sum_{k=2}^{n} \frac{1}{k} = 2n\left(H_n - 1\right)
    \leq 2n \ln n + O(n),
  \end{align*}
  wobei $H_n = \sum_{k=1}^n \frac{1}{k} \approx \ln n + \gamma$ die $n$-te
  harmonische Zahl ist.
\end{proof}

Die Worst-Case-Laufzeit bleibt $\Theta(n^2)$, tritt aber nur mit exponentiell
kleiner Wahrscheinlichkeit auf: Per Chernoff-Schranke (Satz~\ref{satz:chernoff})
gilt $T(n) > cn^2$ mit Wahrscheinlichkeit $e^{-\Omega(n)}$.

\subsection{Kargerscher Algorithmus für minimale Schnitte}
\label{subsec:karger}

\subsubsection{Das Problem und die Idee}

\begin{definition}[Minimaler Schnitt]\label{def:mincut}
  Gegeben ein ungerichteter, zusammenhängender Multigraph $G=(V,E)$ mit
  $|V|=n$ Knoten. Ein \textbf{Schnitt} ist eine Partition $V = S \cup \bar{S}$;
  seine \textbf{Kapazität} ist $|E(S,\bar{S})|$, die Anzahl der Kanten
  zwischen $S$ und $\bar{S}$. Gesucht ist ein Schnitt minimaler Kapazität.
\end{definition}

Kargers randomisierter Kontraktion salgorithmus \cite{karger1993global}
funktioniert durch wiederholte \emph{Kantenkontraktion}:

\begin{algorithm}[H]
  \caption{Karger-Kontraktion}
  \label{alg:karger}
  \begin{algorithmic}[1]
    \Require{Multigraph $G = (V, E)$ mit $n = |V| \geq 2$}
    \Ensure{Kandidat für einen minimalen Schnitt}
    \State $G' \gets G$
    \While{$|V(G')| > 2$}
      \State Wähle eine Kante $e = \{u,v\} \in E(G')$ uniformly at random.
      \State Kontrahiere $e$: Verschmelze $u$ und $v$ zu einem Superknoten.
        \Comment{Schleifen werden entfernt, Mehrfachkanten bleiben}
      \State Entferne alle Schleifen aus $G'$.
    \EndWhile
    \State \Return die Menge $E(G')$ (alle verbleibenden Kanten bilden
           einen Kandidat-Schnitt)
  \end{algorithmic}
\end{algorithm}

\subsubsection{Analyse: Wahrscheinlichkeit des Erfolgs}

\begin{satz}[Erfolgswahrscheinlichkeit der Karger-Kontraktion]\label{satz:kargerprob}
  Sei $k$ die Größe eines minimalen Schnitts. Der Algorithmus gibt diesen
  minimalen Schnitt mit Wahrscheinlichkeit $\geq \frac{2}{n(n-1)}$ zurück.
\end{satz}

\begin{proof}
  Sei $F$ die Kantenmenge eines festen minimalen Schnitts mit $|F| = k$.
  Da jeder Knoten mindestens $k$ Kanten hat, gilt $|E| \geq \frac{nk}{2}$.

  Beim $i$-ten Schritt enthält der Graph noch $n-i+1$ Superknoten, also
  mindestens $\frac{(n-i+1)k}{2}$ Kanten. Die Wahrscheinlichkeit, eine Kante
  aus $F$ nicht zu kontrahieren, ist mindestens
  \[
    1 - \frac{k}{\frac{(n-i+1)k}{2}} = 1 - \frac{2}{n-i+1}.
  \]
  Die Wahrscheinlichkeit, dass \emph{keine} der $n-2$ Kontraktionsschritte
  eine Kante aus $F$ trifft, ist mindestens:
  \begin{align*}
    \prod_{i=1}^{n-2}\!\left(1 - \frac{2}{n-i+1}\right)
    &= \prod_{t=3}^{n} \frac{t-2}{t}
    = \frac{1 \cdot 2 \cdot 3 \cdots (n-2)}{3 \cdot 4 \cdots n}
    = \frac{2}{n(n-1)}.
  \end{align*}
  Hierbei wurde Vollständigkeit genutzt: Die Wahrscheinlichkeiten der $k$
  Kantenwahlen summieren sich in jedem Schritt zu~$1$.
\end{proof}

\subsubsection{Karger-Stein-Algorithmus und Gesamtlaufzeit}

Durch $\Theta(n^2 \log n)$-fache Wiederholung und rekursive Beschleunigung
\cite{karger1996new} erhält man den \textbf{Karger-Stein-Algorithmus} mit:
\begin{itemize}
  \item Erfolgswahrscheinlichkeit $\geq 1 - \frac{1}{n}$ (hohe Wahrscheinlichkeit).
  \item Laufzeit $O(n^2 \log^3 n)$.
\end{itemize}

\subsubsection{Laufzeit einer Einzeliteration}

Eine Kontraktion benötigt $O(n)$ Zeit (Aktualisierung der Adjazenzlisten).
Es gibt $n-2$ Kontraktionen, daher ist eine Iteration in $O(n^2)$ durchführbar.

\subsection{Miller-Rabin-Primzahltest}
\label{subsec:miller}

\subsubsection{Motivation}

Der deterministische Primzahltest hat lange Zeit Schwierigkeiten bereitet.
AKS (2002) liefert einen deterministischen Polynomialzeitalgorithmus
\cite{agrawal2004primes}, ist aber in der Praxis langsamer als der
randomisierte Miller-Rabin-Test.

\subsubsection{Mathematische Grundlagen}

\begin{definition}[Zusammengesetzte Zahl und Fermat-Test]\label{def:fermat}
  Für eine Primzahl $p$ gilt nach dem Kleinen Fermatschen Satz:
  $a^{p-1} \equiv 1 \pmod{p}$ für alle $a$ mit $\gcd(a,p)=1$.
  Schreibe $p - 1 = 2^s \cdot d$ mit $d$ ungerade. Dann gilt für alle $a$:
  Entweder $a^d \equiv 1 \pmod{p}$, oder es existiert ein $r \in \{0,\ldots,s-1\}$
  mit $a^{2^r d} \equiv -1 \pmod{p}$.
  Trifft keines zu, ist $p$ zusammengesetzt.
\end{definition}

\begin{algorithm}[H]
  \caption{Miller-Rabin-Primzahltest}
  \label{alg:miller}
  \begin{algorithmic}[1]
    \Require{Ungerade Zahl $n \geq 3$; Sicherheitsparameter $t \geq 1$}
    \Ensure{\texttt{ZUSAMMENGESETZT} oder \texttt{WAHRSCHEINLICH PRIM}}
    \State Schreibe $n - 1 = 2^s \cdot d$ mit $d$ ungerade.
    \For{$i \gets 1$ \textbf{to} $t$}
      \State Wähle $a \in \{2, 3, \ldots, n-2\}$ gleichverteilt zufällig.
      \State $x \gets a^d \bmod n$ \Comment{Schnelle Exponentiation: $O(\log n)$ Multiplikationen}
      \If{$x = 1$ \textbf{oder} $x = n-1$}
        \State \textbf{continue}
      \EndIf
      \State $\mathit{found} \gets \texttt{false}$
      \For{$r \gets 1$ \textbf{to} $s-1$}
        \State $x \gets x^2 \bmod n$
        \If{$x = n-1$}
          \State $\mathit{found} \gets \texttt{true}$; \textbf{break}
        \EndIf
      \EndFor
      \If{\textbf{not} $\mathit{found}$}
        \State \Return \texttt{ZUSAMMENGESETZT}
      \EndIf
    \EndFor
    \State \Return \texttt{WAHRSCHEINLICH PRIM}
  \end{algorithmic}
\end{algorithm}

\subsubsection{Fehlerwahrscheinlichkeit und Laufzeit}

\begin{satz}[Fehlerwahrscheinlichkeit des Miller-Rabin-Tests]\label{satz:mrfehler}
  Ist $n$ zusammengesetzt, so gilt: Für eine gleichverteilt zufällige
  Basis $a \in \{2, \ldots, n-2\}$ ist
  \[
    \Prob\!\left(\text{$a$ ein „starker Lügner" für $n$}\right)
    \leq \frac{1}{4}.
  \]
  Nach $t$ unabhängigen Iterationen:
  \[
    \Prob\!\left(\text{$n$ zusammengesetzt und Algorithmus sagt PRIM}\right)
    \leq \left(\frac{1}{4}\right)^{\!t}.
  \]
\end{satz}

Mit $t = 40$ gilt die Fehlerwahrscheinlichkeit $\leq 4^{-40} \approx 10^{-24}$,
also kleiner als die Fehlerwahrscheinlichkeit moderner Hardware.

\textbf{Vollständigkeit:} In jeder Iteration gilt
$\Prob(\text{kein Lügner}) + \Prob(\text{Lügner}) = 1$.

\textbf{Laufzeit:}
\begin{itemize}
  \item Exponentiation $a^d \bmod n$: $O(\log n)$ modulare Multiplikationen,
    jede $O(\log^2 n)$ bitweise, also $O(\log^3 n)$ pro Basis.
  \item Quadrierungen: $s \leq \log_2 n$ Schritte, je $O(\log^2 n)$.
  \item Gesamtlaufzeit: $O(t \log^3 n)$.
  \item Mit FFT-basierter Multiplikation: $O(t \log^2 n \cdot \log\log n)$.
\end{itemize}

\subsection{Polynomial Identity Testing (Schwartz-Zippel-Lemma)}
\label{subsec:schwartz}

\subsubsection{Polynomial Identity Testing}

\begin{definition}[Polynomial Identity Testing (PIT)]\label{def:pit}
  Gegeben zwei Polynome $P, Q \in \mathbb{F}[x_1, \ldots, x_n]$ mit Grad
  $\leq d$, jeweils als Schaltkreis gegeben. Entscheide, ob $P \equiv Q$.
\end{definition}

\begin{satz}[Schwartz-Zippel-Lemma]\label{satz:sz}
  Sei $\mathbb{F}$ ein Körper, $S \subseteq \mathbb{F}$ eine endliche
  nicht-leere Teilmenge und $P \in \mathbb{F}[x_1,\ldots,x_n]\setminus\{0\}$
  ein Polynom vom Grad $\leq d$. Wählt man $r_1, \ldots, r_n$ unabhängig
  gleichverteilt aus $S$, so gilt:
  \[
    \Prob(P(r_1, \ldots, r_n) = 0) \leq \frac{d}{|S|}.
  \]
\end{satz}

\begin{algorithm}[H]
  \caption{Randomisierter Polynom-Identitätstest (Schwartz-Zippel)}
  \label{alg:sz}
  \begin{algorithmic}[1]
    \Require{Polynome $P, Q$ (als Schaltkreise), endliche Menge $S \subseteq
             \mathbb{F}$ mit $|S| \geq 2d$, Grad $d$}
    \Ensure{\texttt{IDENTISCH} oder \texttt{VERSCHIEDEN}}
    \State $R \gets P - Q$
    \For{$i \gets 1$ \textbf{to} $n$}
      \State Wähle $r_i \in S$ gleichverteilt und unabhängig.
    \EndFor
    \If{$R(r_1, \ldots, r_n) = 0$}
      \State \Return \texttt{IDENTISCH} \Comment{Falsch-Positiv mit $\Prob \leq d/|S| \leq 1/2$}
    \Else
      \State \Return \texttt{VERSCHIEDEN} \Comment{Immer korrekt}
    \EndIf
  \end{algorithmic}
\end{algorithm}

\textbf{Laufzeit:} $O(n \cdot d \cdot \text{Schaltkreisgröße})$ für eine
Auswertung. Mit $|S| = 2d$ ist Fehlerwahrscheinlichkeit $\leq \frac{1}{2}$;
durch $k$-fache Wiederholung auf $2^{-k}$ reduzierbar.

\subsection{Hashing und Universal Hashing}
\label{subsec:hashing}

\subsubsection{Motivation}

Hash-Tabellen erreichen $O(1)$ amortisierte Laufzeit für Einfügen, Löschen
und Suchen – \emph{im Erwartungswert} unter randomisiertem Hashing.

\begin{definition}[Universelle Hashfamilie]\label{def:uhash}
  Eine Familie $\mathcal{H}$ von Hashfunktionen $h: U \to \{0,\ldots,m-1\}$
  heißt \textbf{universell}, wenn für alle $x \neq y \in U$:
  \[
    \Prob_{h\sim\mathcal{H}}(h(x) = h(y)) \leq \frac{1}{m}.
  \]
\end{definition}

\begin{algorithm}[H]
  \caption{Randomisiertes Hashing (Chaining)}
  \label{alg:hashing}
  \begin{algorithmic}[1]
    \Require{Universum $U$, Tabellengröße $m$, universelle Hashfamilie $\mathcal{H}$}
    \Ensure{Hash-Tabelle mit erwartet $O(1)$ Operationen}
    \State Wähle $h \in \mathcal{H}$ gleichverteilt.
    \State Initialisiere Tabelle $T[0..m-1]$ mit leeren Listen.
    \Statex
    \Function{Insert}{$x$}
      \State Füge $x$ an $T[h(x)]$ an.
    \EndFunction
    \Statex
    \Function{Search}{$x$}
      \State Durchsuche Liste $T[h(x)]$ nach $x$.
      \State \Return \texttt{GEFUNDEN} oder \texttt{NICHT GEFUNDEN}
    \EndFunction
  \end{algorithmic}
\end{algorithm}

\subsubsection{Erwartete Listenlänge}

\begin{satz}[Erwartete Listenlänge bei universalem Hashing]\label{satz:hashew}
  Sei $S \subseteq U$ die gespeicherte Menge mit $|S| = n$ und $m = n$.
  Für jedes $x \in S$ gilt:
  \[
    \E\!\left[\text{Anzahl Kollisionen mit }x\right]
    = \sum_{y \in S,\, y \neq x} \Prob(h(x) = h(y))
    \leq (n-1) \cdot \frac{1}{n} < 1.
  \]
  Daher ist die erwartete Suchzeit für jedes Element $O(1)$.
\end{satz}

\subsection{Randomisierter Las-Vegas-Algorithmus: Randomisiertes Finden des Minimums}
\label{subsec:randmin}

\begin{algorithm}[H]
  \caption{Randomisiertes Minimum-Finden (LazySelect)}
  \label{alg:lazyselect}
  \begin{algorithmic}[1]
    \Require{Array $A[1..n]$, Position $k$ (gesucht: $k$-kleinstes Element)}
    \Ensure{Das $k$-kleinste Element von $A$}
    \State $r \gets n^{3/4}$
    \State Wähle eine Stichprobe $R$ von $r$ Elementen aus $A$ gleichverteilt
           mit Zurücklegen.
    \State Sortiere $R$.
    \State $\ell \gets \max\!\left(1,\, \lfloor kn^{-1/4} \rfloor - \sqrt{n}\right)$,
           $\;h \gets \min\!\left(r,\, \lceil kn^{-1/4} \rceil + \sqrt{n}\right)$
    \State $a \gets R[\ell]$, $b \gets R[h]$
    \State $P \gets \{x \in A : a \leq x \leq b\}$
    \If{$k$-kleinstes Element von $A$ liegt in $P$ \textbf{und} $|P| \leq 4n^{3/4}+2$}
      \State Sortiere $P$ und gib das $(k - |\{x \in A : x < a\}|)$-te Element zurück.
    \Else
      \State Starte neu (Misserfolg mit Wahrscheinlichkeit $O(n^{-1/4})$).
    \EndIf
  \end{algorithmic}
\end{algorithm}

\textbf{Laufzeit:} Erwartet $O(n)$ Vergleiche – optimal für das Selektionsproblem.

% ===============================================================
\section{Die probabilistische Methode}
\label{sec:probmethode}
% ===============================================================

\subsection{Grundprinzip}

Die \emph{probabilistische Methode} \cite{alon2016probabilistic} ist eine
nicht-konstruktive Beweistechnik: Man zeigt, dass ein bestimmtes Objekt
(Graph, Belegung, Kode) existiert, indem man zeigt, dass ein zufällig
gewähltes Objekt die gewünschte Eigenschaft mit positiver Wahrscheinlichkeit
besitzt.

\paragraph*{Probabilistische Methode – Grundschema.}
\begin{quote}
  Sei $\mathcal{O}$ eine Menge von Objekten und $\mathcal{E}$ eine Eigenschaft.
  Man definiert eine Wahrscheinlichkeitsverteilung auf $\mathcal{O}$ und zeigt:
  \[
    \Prob(O \in \mathcal{O} \text{ hat Eigenschaft } \mathcal{E}) > 0.
  \]
  Daraus folgt: Es existiert ein Objekt in $\mathcal{O}$ mit Eigenschaft $\mathcal{E}$.
\end{quote}

\subsection{Beispiel: Turán-Ramsey und unabhängige Mengen}

\begin{satz}[Ramsey-Schranke via probabilistische Methode]\label{satz:ramsey}
  Es gilt $R(k,k) \geq 2^{k/2}$ für alle $k \geq 3$.
\end{satz}

\begin{proof}
  Betrachte $K_n$ mit $n = \lfloor 2^{k/2} \rfloor$. Färbe jede Kante
  unabhängig mit Wahrscheinlichkeit $\frac{1}{2}$ rot oder blau. Für eine
  feste $k$-Knotenmenge $S$ gilt:
  \[
    \Prob(\text{alle Kanten in }S\text{ einfarbig}) = 2 \cdot 2^{-\binom{k}{2}}.
  \]
  Mit der Union Bound (Satz~\ref{satz:grundregeln}):
  \[
    \Prob(\exists \text{ einfarbige } k\text{-Clique}) \leq \binom{n}{k} \cdot 2^{1-\binom{k}{2}}.
  \]
  Für $n < 2^{k/2}$ ergibt sich, dass diese Schranke $< 1$ ist, also mit
  positiver Wahrscheinlichkeit existiert keine einfarbige $k$-Clique –
  und daher existiert eine Färbung ohne solche.
\end{proof}

\subsection{Lovász Local Lemma}

\begin{satz}[Lovász Local Lemma (symmetrische Form)]\label{satz:lll}
  Seien $A_1, \ldots, A_n$ Ereignisse mit $\Prob(A_i) \leq p$ für alle $i$.
  Jedes $A_i$ hängt von höchstens $d$ anderen ab. Falls $ep(d+1) \leq 1$,
  dann gilt:
  \[
    \Prob\!\left(\bigcap_{i=1}^n \overline{A_i}\right) > 0.
  \]
\end{satz}

Das LLL ist eine der tiefsten Anwendungen der probabilistischen Methode und
hat weitreichende algorithmische Konsequenzen (Konstruktiver Beweis via
Moser-Tardos \cite{moser2010constructive}).

% ===============================================================
\section{Gesetz großer Zahlen und Zentraler Grenzwertsatz}
\label{sec:ggz}
% ===============================================================

\subsection{Gesetz der großen Zahlen}

\begin{satz}[Schwaches Gesetz der großen Zahlen (GGZ)]\label{satz:sggz}
  Seien $X_1, X_2, \ldots$ i.i.d.\ ZV mit $\E[X_i] = \mu < \infty$. Dann gilt
  für alle $\varepsilon > 0$:
  \[
    \Prob\!\left(\left|\frac{X_1 + \cdots + X_n}{n} - \mu\right| > \varepsilon\right) \to 0
    \quad \text{für } n \to \infty.
  \]
\end{satz}

Das GGZ rechtfertigt die Interpretation des Erwartungswerts als
\emph{langfristiger Durchschnitt}: Wenn man das Experiment unendlich oft
wiederholt und die Ergebnisse summiert, summieren sich die
Wahrscheinlichkeiten aller Experimente zu $n \cdot 1 = n$, und der
Durchschnitt konvergiert zu $\mu$.

\subsection{Zentraler Grenzwertsatz}

\begin{satz}[Zentraler Grenzwertsatz (Lindeberg-Lévy, ZGS)]\label{satz:zgs}
  Seien $X_1, X_2, \ldots$ i.i.d.\ ZV mit $\E[X_i] = \mu$ und
  $\Var(X_i) = \sigma^2 \in (0,\infty)$. Dann gilt für alle $x \in \mathbb{R}$:
  \[
    \lim_{n\to\infty} \Prob\!\left(\frac{X_1+\cdots+X_n - n\mu}{\sigma\sqrt{n}} \leq x\right)
    = \Phi(x) = \frac{1}{\sqrt{2\pi}}\int_{-\infty}^x e^{-t^2/2}\,dt.
  \]
\end{satz}

Der ZGS erklärt die Allgegenwart der Normalverteilung in der Natur und bildet die Grundlage der statistischen Inferenz. Dass der ZGS aussagt, dass die Normalverteilung die Verteilung für Beschreibungen in der Natur ist, ist nicht überraschend. Lindeberg-Lévy hat dies richtig beschrieben und das muss auch so sein, dass wir erst über die Betrachtung der Normalverteilung und ihrer Notwendigkeit in der Natur klar sehen, dass es eben doch bewusste Absichten in den scheinbar zufälligen Ereignissen der Natur gibt.

% ===============================================================
\section{Schluss}
\label{sec:schluss}
% ===============================================================

In dieser Arbeit wurde die Lehre vom Zufall in der Mathematik entwickelt.
Ausgangspunkt war das einfache Beispiel des Würfelns eines Würfels. An diesem
Beispiel wurden die grundlegenden Begriffe eingeführt: der Versuch, die
Ereignisse, die Menge aller Ereignisse $\Omega$, und die Wahrscheinlichkeit
für den Zufall eines Ereignisses. Es wurde gezeigt, dass ein Versuch
mathematisch vollständig beschrieben ist, wenn die Summe aller
Wahrscheinlichkeiten der Zufälle aller Ereignisse insgesamt $1$ ergibt.

Für die Analyse von Versuchen wurde das Konzept der Zufallsvariable
eingeführt. Eine Zufallsvariable beschreibt den Zufall für das Ereignis
variabel für den konkreten Ausgang. Mit Hilfe der Zufallsvariable kann der
Erwartungswert berechnet werden. Der Erwartungswert beschreibt formal am
genauesten, welcher Ausgang bei einem Versuch zu erwarten ist. Es wurde aber
auch gewarnt, dass der Erwartungswert in der Praxis nicht immer sinnvoll
interpretierbar ist. Beim Würfeln eines Würfels ergibt sich ein
Erwartungswert von $3{,}5$, eine Augenzahl, die es nicht gibt. Der
Erwartungswert beschreibt das, was im Mittel über viele Wiederholungen des
Versuchs zu erwarten ist.

Am Beispiel des 3-SAT-Algorithmus mit Zufall wurde gezeigt, wie Zufall in
der Informatik eingesetzt wird, um effiziente Algorithmen zu entwickeln. Der
Algorithmus erreicht eine Approximationsgüte von $\frac{7}{8}$: Im
Erwartungswert werden mindestens $87{,}5\,\%$ aller erfüllbaren Klauseln
erfüllt. Die mathematische Analyse zeigt, dass
$\E[Y] = \frac{7}{8}m \ge \frac{7}{8}\,O$, wobei $O$ die optimale Anzahl
erfüllbarer Klauseln ist. Dies ist ein Beispiel dafür, dass die
mathematische Beschreibung des Zufalls in der Praxis anwendbar ist und zu
verlässlichen Ergebnissen führt.

Die Wahrscheinlichkeitsverteilungen sind zentrale Werkzeuge für die Analyse
von Versuchen mit Zufall. Es wurden mehrere Verteilungen vorgestellt: die
Bernoulli-Verteilung für Versuche mit zwei möglichen Ereignissen, die
Binomialverteilung für die Anzahl der Erfolge bei $n$ Versuchen, die
geometrische Verteilung für die Anzahl der Versuche bis zum ersten Erfolg,
sowie die Gaußsche Verteilung für kontinuierliche Ereignisse. Jede dieser
Verteilungen hat ihre spezifischen Anwendungsbereiche. Es wurde mehrfach
betont, dass die Wahl der richtigen Verteilung entscheidend ist: Kleinste
Details in der Beschreibung des Versuchs können darüber entscheiden, welche
Verteilung anwendbar ist. Wird eine andere Verteilung gewählt, führt dies
zu anderen Ergebnissen in der Analyse.

Die Lehre vom Zufall in der Mathematik muss mit Genauigkeit und Vorsicht
entwickelt und angewendet werden. Wann immer die mathematische Beschreibung
so weit von der Praxis entfernt ist, dass sie nicht mehr anwendbar ist, wird
das eigentliche Ziel der Mathematik in diesem Thema verletzt. Die Mathematik
soll helfen, den Zufall zu verstehen und in der Praxis verlässliche
Entscheidungen zu treffen.

Das in dieser Arbeit entwickelte Bild zeigt: Die Lehre vom Zufall ist kein
Randgebiet der Mathematik, sondern ein lebendiges Kerngebiet, das in alle
Bereiche der modernen Wissenschaft und Technik hineinwirkt. Im Folgenden
werden aktuelle Forschungsrichtungen und offene Fragen ausführlich diskutiert.

\subsection{Komplexitätstheorie und derandomisierte Algorithmen}

Eine der tiefsten offenen Fragen der Theoretischen Informatik lautet:

\paragraph*{Offene Frage: P vs.\ BPP.}
\begin{quote}
  Gilt $\text{P} = \text{BPP}$? Ist jede Sprache, die ein
  probabilistischer Polynomialzeit-Algorithmus mit Fehlerwahrscheinlichkeit
  $\leq 1/3$ entscheidet, auch deterministisch in Polynomialzeit entscheidbar?
\end{quote}

Impagliazzo und Wigderson \cite{impagliazzo1997p} zeigten, dass falls
exponentielle untere Schranken für bestimmte Schaltkreisfamilien gelten
(konkret: falls $\text{EXP} \not\subseteq \text{P/poly}$), dann ist
$\text{P} = \text{BPP}$. Dies weist den Weg zu einer vollständigen
Derandomisierung: \emph{Pseudozufallsgeneratoren} von ausreichender Güte
würden den Zufall in Polynomialzeitalgorithmen überflüssig machen. Aktuelle
Forschung (z.\,B. Ta-Shma \cite{tashma2017explicit}, Braverman et al.)
arbeitet an optimalen expliziten Konstruktionen von
Pseudozufallsgeneratoren mit minimalem Seed-Length.

\subsection{Quantencomputing und Quanten-Wahr\-schein\-lich\-keits\-theo\-rie}

Quantenalgorithmen erweitern das Modell des randomisierten Algorithmus
fundamental: Statt klassischer Wahrscheinlichkeitsamplituden arbeiten
Quantenalgorithmen mit \emph{komplexen Amplituden} $\alpha_i \in \mathbb{C}$,
wobei die Normierungsbedingung lautet:
\[
  \sum_i |\alpha_i|^2 = 1.
\]
Dies ist die Quantenversion des Vollständigkeitsprinzips. Interference zwischen
Pfaden ermöglicht \emph{exponentielle Beschleunigung} für bestimmte Probleme:
\begin{itemize}
  \item \textbf{Shors Algorithmus} \cite{shor1994algorithms}: Primfaktorzerlegung
    in $O((\log n)^3)$ Quantengattern – gegenüber sub-exponetiell im Klassischen.
  \item \textbf{Grovers Suchalgorithmus} \cite{grover1996fast}: Datenbanksuche
    in $O(\sqrt{N})$ statt $O(N)$ klassisch. Dieser quadratische Speedup ist
    optimal \cite{bennett1997strengths}.
  \item \textbf{Quantum Approximate Optimization Algorithm (QAOA)}:
    Hybridalgorithmus für kombinatorische Optimierung, der auf randomisierten
    Variationsprinzipien beruht \cite{farhi2014quantum}.
\end{itemize}
Offene Fragen betreffen die genaue Abgrenzung von $\text{BQP}$
(Bounded-Error Quantum Polynomial Time) gegenüber $\text{BPP}$ und $\text{NP}$.
Es ist unbekannt, ob $\text{NP} \subseteq \text{BQP}$.

\subsection{Maschinelles Lernen und Stochastische Optimierung}

Randomisierung ist das Herzstück moderner lerntheoretischer Verfahren:

\subsubsection{Stochastischer Gradientenabstieg (SGD)}
Der \emph{Stochastic Gradient Descent} \cite{robbins1951stochastic} approximiert
den echten Gradienten $\nabla L(\theta) = \frac{1}{N}\sum_i \nabla \ell_i(\theta)$
durch einen zufälligen Mini-Batch-Gradienten:
\[
  \theta_{t+1} = \theta_t - \eta_t \cdot \frac{1}{|B_t|} \sum_{i \in B_t}
  \nabla \ell_i(\theta_t),
\]
wobei $B_t \subseteq \{1,\ldots,N\}$ ein zufällig gezogener Mini-Batch ist.
Das Vollständigkeitsprinzip in diesem Kontext: Die Ziehung des Mini-Batches
ist eine Wahrscheinlichkeitsverteilung auf $\binom{N}{b}$ Teilmengen der
Größe $b$, die zu $1$ normiert ist.

\subsubsection{Dropout und Regularisierung via Zufall}
\emph{Dropout} \cite{srivastava2014dropout} deaktiviert während des Trainings
jedes Neuron unabhängig mit Wahrscheinlichkeit $p$: Ein randomisiertes
Ensemble von $2^n$ dünnen Netzwerken wird implizit trainiert. Die
Vollständigkeitsbedingung – $p + (1-p) = 1$ – sichert, dass jedes Neuron
entweder aktiv oder inaktiv ist.

\subsubsection{Convolutional Neural Networks und Zufälligkeit}
Zufällige Initialisierung (He-Initialisierung, Xavier-Initialisierung)
und Data Augmentation (zufällige Spiegelungen, Ausschnitte, Farbveränderungen)
sind essentiell für den Trainingserfolg moderner tiefer Netze.

\subsection{Kryptographie und Informationstheorie}

\subsubsection{Perfekte Sicherheit und Einmalblock-Verschlüsselung}
Shannon zeigte \cite{shannon1949communication}, dass ein kryptographisches
System genau dann \emph{perfekt sicher} ist, wenn
\[
  \Prob(M = m \mid C = c) = \Prob(M = m)
\]
für alle Klartexte $m$ und Chiffretexte $c$ gilt – d.\,h., das Beobachten des
Chiffretextes ändert die Wahrscheinlichkeit des Klartextes nicht. Das
One-Time-Pad erreicht perfekte Sicherheit genau dann, wenn der Schlüssel
gleichverteilt und so lang wie der Klartext ist (Vollständigkeit!).

\subsubsection{Pseudozufallsgeneratoren in der Kryptographie}
Kryptographisch sichere Pseudozufallsgeneratoren (CSPRNG) erzeugen aus
kurzen, echten Zufallsquellen (Seeds) lange, ununterscheidbare Bitfolgen.
Grundlegende Konstruktionen:
\begin{itemize}
  \item AES-CTR-DRBG (NIST-Standard SP 800-90A),
  \item Blum-Blum-Shub-Generator \cite{blum1986simple}, beweisbar sicher
    unter Annahme der Schwierigkeit der Quadratwurzel modulo $n=pq$.
\end{itemize}

\subsubsection{Zero-Knowledge-Beweise}
Interaktive Beweissysteme \cite{goldwasser1989knowledge} nutzen Zufall
fundamental: Der Verifier stellt zufällige Challenges, und der Prover muss
sie korrekt beantworten – mit Wahrscheinlichkeit $1/2$ pro Runde scheitert
ein Betrüger. Die \emph{Vervollständigung} (Completeness) ist: Ein ehrlicher
Prover beantwortet immer korrekt ($\Prob = 1$); die \emph{Solidität}
(Soundness) garantiert, dass ein Betrüger mit $\Prob \leq 1/2$ erwischt wird.

\subsection{Randomisierte Datenstrukturen}

\subsubsection{Skip-Listen}
Pugh \cite{pugh1990skip} entwickelte die Skip-Liste: Eine probabilistische
Alternative zu Suchbäumen mit erwarteter $O(\log n)$ Such-, Einfüge- und
Löschzeit. Jedes Element wird auf Level $k$ gefördert mit Wahrscheinlichkeit
$p^k$ – eine geometrische Verteilung, deren Vollständigkeit
$\sum_{k=0}^{\infty} p^k(1-p) = 1$ garantiert, dass jedes Element auf
einem endlichen Level endet.

\subsubsection{Bloom-Filter}
Bloom \cite{bloom1970space} führte eine randomisierte Mengendarstellung
mit Falsch-Positiv-Rate $\varepsilon$ und ohne Falsch-Negativ ein:
\[
  \Prob(\text{false positive}) \approx \left(1 - e^{-kn/m}\right)^k,
\]
wobei $n$ die Elementanzahl, $m$ die Bitgröße des Filters und $k$ die Zahl
der Hashfunktionen ist. Optimal ist $k = \frac{m}{n}\ln 2$.

\subsubsection{Treaps}
Ein \emph{Treap} \cite{seidel1996randomized} kombiniert einen binären
Suchbaum mit einem Max-Heap, wobei die Heap-Prioritäten gleichverteilt
zufällig zugewiesen werden. Die erwartete Suchtiefe beträgt $O(\log n)$,
und die Struktur verhält sich wie ein zufällig aufgebauter Suchbaum.

\subsection{Approximationsalgorithmen und semidefinite Programmierung}

\subsubsection{Goemans-Williamson-Algorithmus für MAX-CUT}
Goemans und Williamson \cite{goemans1995improved} zeigten mittels
\emph{Semidefiniter Programmierung} (SDP) und anschließendem
Runden durch einen zufälligen Hyperebenen-Schnitt, dass MAX-CUT mit
Faktor $\alpha_{\text{GW}} = \frac{2}{\pi}\int_0^\pi \frac{1-\cos\theta}{2}
\cdot \frac{d\theta}{\pi} \approx 0{,}878$ approximierbar ist.

Das Runden: Sei $v$ ein gleichverteilter zufälliger Einheitsvektor in $\mathbb{R}^n$.
Weise Knoten $i$ der Seite $+1$ zu, falls $v \cdot x_i \geq 0$, und $-1$ sonst.
Die Vollständigkeit: $\Prob(v \cdot x_i \geq 0) + \Prob(v \cdot x_i < 0) = 1$.

\subsubsection{Håstads 3-SAT-Härte und optimale Approximierbarkeit}
Håstad \cite{hastad2001some} bewies, dass für E3-SAT kein Approximationsalgorithmus
mit Faktor $> 7/8$ existiert (unter $\text{P}\neq\text{NP}$), womit der
in Abschnitt~\ref{subsec:3sat} beschriebene einfache Algorithmus optimal ist.

\subsection{Monte-Carlo-Simulation und numerische Integration}

\subsubsection{Monte-Carlo-Integration}
Für die numerische Integration einer Funktion $f: [0,1]^d \to \mathbb{R}$:
\[
  I = \int_{[0,1]^d} f(x)\,dx \approx \hat{I}_n = \frac{1}{n}\sum_{i=1}^n f(U_i),
\]
wobei $U_1, \ldots, U_n$ gleichverteilt zufällig in $[0,1]^d$ gewählt werden.
Das GGZ garantiert $\hat{I}_n \to I$, der ZGS liefert die Konvergenzrate:
\[
  \Var(\hat{I}_n) = \frac{\sigma^2_f}{n}, \quad
  \text{Standardfehler} = \frac{\sigma_f}{\sqrt{n}}.
\]
Im Unterschied zu deterministischen Quadraturregeln, die für $d$-dimensional
unter dem \emph{Fluch der Dimensionalität} leiden, skaliert der
Monte-Carlo-Fehler unabhängig von $d$ mit $O(n^{-1/2})$.

\subsubsection{Markov-Ketten-Monte-Carlo (MCMC)}
\emph{Markov-Ketten-Monte-Carlo}-Methoden wie der \textbf{Metropolis-Hastings-Algorithmus}
\cite{metropolis1953equation, hastings1970monte} und
\textbf{Gibbs-Sampling} erzeugen Stichproben aus komplexen
Wahrscheinlichkeitsverteilungen $\pi$ durch eine Markov-Kette, die $\pi$ als
stationäre Verteilung hat. Ihre Vollständigkeit:
$\sum_y \Prob(X_{t+1}=y \mid X_t=x) = 1$ (die Übergangsprobabilitäten bilden
eine vollständig beschriebene Verteilung für jeden Zustand $x$).

\subsubsection{Simulated Annealing}
\emph{Simulated Annealing} \cite{kirkpatrick1983optimization} ist ein
randomisiertes metaheuristisches Optimierungsverfahren: Schlechtere
Lösungen werden mit Wahrscheinlichkeit
$\Prob(\text{Akzeptanz}) = e^{-\Delta E / T}$
akzeptiert, wobei $T$ eine fallende „Temperatur" ist. Die Normierung
$\Prob(\text{Akzeptanz}) + \Prob(\text{Ablehnung}) = 1$ stellt die
vollständige Beschreibbarkeit des Übergangsereignisses sicher.

\subsection{Streaming und Sublinear-Zeit-Algorithmen}

\subsubsection{Frequenz-Schätzung und Count-Min-Sketch}
Cormode und Muthukrishnan \cite{cormode2005improved} entwickelten
Count-Min-Sketch zur Approximation von Element-Häufigkeiten in Datenströmen:
Mit $O(\frac{1}{\varepsilon} \log \frac{1}{\delta})$ Speicherplatz wird die
Häufigkeit jedes Elements bis auf $\varepsilon n$ mit Wahrscheinlichkeit
$1-\delta$ korrekt ausgegeben. Vollständigkeit:
$\Prob(\text{Fehler} \leq \varepsilon n) + \Prob(\text{Fehler} > \varepsilon n)= 1$.

\subsubsection{$F_0$-Schätzung (distinct elements)}
Flajolet und Martin \cite{flajolet1985probabilistic} entwickelten den ersten
probabilistischen Algorithmus zur Schätzung der Anzahl verschiedener Elemente
in einem Datenstrom – ein Problem, das deterministisch $\Omega(n)$ Speicher
erfordert, randomisiert aber in $O(\log n)$ Bits gelöst werden kann.

\subsection{Randomisierung in der Bioinformatik und Computational Biology}

\subsubsection{BLAST und Heuristischer Sequenzvergleich}
Der BLAST-Algorithmus \cite{altschul1990basic} nutzt randomisierte Heuristiken
für den lokalem Sequenzvergleich: Durch zufälliges Sampling von Keimwörtern
und statistische Modellierung der Erwartungswerte von Treffern (E-Value)
unter der Annahme zufälliger Sequenzen wird die Sensitivität/Spezifität-Abwägung
quantifiziert. Das E-Value ist gerade der Erwartungswert der Anzahl Treffer
in einer Zufallsdatenbank.

\subsubsection{Phylogenetische Rekonstruktion via Bootstrap}
Felsenstein \cite{felsenstein1985confidence} führte das phylogenetische
\emph{Bootstrap-Verfahren} ein: Durch $B$ zufällige Stichproben mit Zurücklegen
aus den Alignment-Spalten werden $B$ Bootstrap-Bäume erstellt; der
Bootstrap-Unterstützungswert eines Astes ist sein Anteil in diesen Bäumen.
Grundlage ist das Vollständigkeitsprinzip: Jede Bootstrap-Stichprobe ist
vollständig beschrieben (alle $n$ Spalten werden gezogen, die Summe der
Stichproben-Wahrscheinlichkeiten ist $1$).

\subsection{Lerntheorie und PAC-Lernen}

\subsubsection{Probably Approximately Correct (PAC) Lernen}
Valiants PAC-Lernmodell \cite{valiant1984theory} quantifiziert, wie viele
Stichproben ein Lernalgorithmus braucht, um mit Wahrscheinlichkeit $\geq 1-\delta$
einen Klassifizierer mit Fehler $\leq \varepsilon$ zu lernen. Die
Stichprobenkomplexität für eine Hypothesenklasse $\mathcal{H}$ mit
VC-Dimension $d$ beträgt:
\[
  m = O\!\left(\frac{d + \log(1/\delta)}{\varepsilon}\right).
\]
Vollständigkeit: Der Lernfehler ist eine ZV; sein Erwartungswert
$\sum_{\varepsilon'} \varepsilon' \Prob(\text{Fehler}=\varepsilon')$ existiert
und ist klein.

\subsection{Offene Probleme und Forschungsfrontieren}

\begin{enumerate}[label=\textbf{\arabic*.}]
  \item \textbf{P vs.\ NP:} Die wichtigste offene Frage der Informatik
    berührt die Lehre vom Zufall: Falls $\text{P} \neq \text{NP}$, existieren
    Probleme, für die randomisierte Algorithmen exponentiell besser als
    deterministische sein könnten \cite{cook1971complexity}.
  \item \textbf{P vs.\ BPP:} Wie oben diskutiert: Ist Randomisierung
    prinzipiell eliminierbar \cite{impagliazzo1997p}?
  \item \textbf{Derandomisierung von $\text{RL}$:} Gilt
    $\text{RL} = \text{L}$? Reingold \cite{reingold2008undirected} zeigte
    dies für ungerichtetes Grapherreichbarkeits-problem; der gerichtete Fall
    ist offen.
  \item \textbf{Optimale Pseudozufallsgeneratoren:} Existieren explizite
    PRGs mit Seed-Länge $O(\log n)$, die alle Polynomialzeit-Tests täuschen?
    Dies würde $\text{P} = \text{BPP}$ implizieren.
  \item \textbf{3-SAT und Random 3-SAT:} Zufällige 3-SAT-Instanzen nahe
    dem Satisfiability-Schwellenwert $\approx 4{,}267$ sind besonders schwer
    \cite{mertens2006threshold}. Das Verhalten randomisierter Algorithmen
    in dieser Phase ist noch nicht vollständig verstanden.
  \item \textbf{MCMC-Mischzeiten:} Für viele natürliche Markov-Ketten
    sind scharfe Schranken an die Mischzeit offen – insbesondere für
    Glauber-Dynamik auf Ising-Modellen bei kritischer Temperatur.
  \item \textbf{Quantum Advantage:} Für welche Problemklassen können
    Quantencomputer exponentiell schneller sein als klassische
    (randomisierte) Computer?
  \item \textbf{Robustheit und Adversarial Machine Learning:} Wie können
    randomisierte Abwehrstrategien (z.\,B.\ Randomized Smoothing
    \cite{cohen2019certified}) robuste neuronale Netze garantieren?
  \item \textbf{Differentially Private Algorithmen:} Randomisierung ist
    zentral für differential privacy \cite{dwork2014algorithmic}:
    Die Ausgabe soll für ähnliche Datenbanken ähnliche Wahrscheinlichkeitsverteilungen
    haben – ein direkter Bezug zum Vollständigkeitsprinzip.
  \item \textbf{Oblivious Algorithmen und Pandemie-Modellierung:}
    Stochastische Epidemiemodelle (SIR, SEIR) \cite{anderson1991infectious}
    nutzen die Lehre vom Zufall, um Ausbruchsdynamiken zu beschreiben;
    Verbesserungen in der Modellierungsgenauigkeit sind von hoher
    gesellschaftlicher Relevanz.
\end{enumerate}

\subsection{Abschlussbemerkung: Das Vollständigkeitsprinzip als\\ epistemologisches Fundament}

Das in der Einleitung formulierte Prinzip – die Summe aller Wahrscheinlichkeiten
ist $1$ – erscheint zunächst wie eine technische Normierungskonvention.
Tatsächlich hat es tiefe epistemologische Bedeutung: Es besagt, dass unser
Modell der Wirklichkeit \emph{vollständig} ist – dass wir alle möglichen
Ausgänge eines Experiments in Betracht gezogen haben. Kein Ergebnis wird
ignoriert, keines wird doppelt gewichtet.

Diese Vollständigkeit ist die mathematische Entsprechung des Prinzips der
\emph{geschlossenen Welt}: Was nicht modelliert ist, kann in der Analyse
nicht auftreten. In der Praxis bedeutet dies eine ständige Warnung:
Wahrscheinlichkeitsmodelle sind nur so gut wie die Vollständigkeit
ihrer Ergebnisräume. Wenn das Modell wichtige Szenarien vergisst, summieren
sich die Wahrscheinlichkeiten zu weniger als $1$ – ein mathematisch
quantifizierbares Signal für Unvollständigkeit.

Von Kolmogorows Axiomen über randomisierte Algorithmen bis hin zu
Quantencomputern bleibt dieses Prinzip der unverrückbare Anker der
Lehre vom Zufall.

\begin{thebibliography}{99}

\bibitem{agrawal2004primes}
M.\ Agrawal, N.\ Kayal, N.\ Saxena:
\textit{PRIMES is in P.}
Annals of Mathematics, 160(2):781--793, 2004.

\bibitem{alon2016probabilistic}
N.\ Alon, J.\ Spencer:
\textit{The Probabilistic Method.}
4.~Auflage.
Wiley, Hoboken, 2016.

\bibitem{altschul1990basic}
S.\,F.\ Altschul, W.\ Gish, W.\ Miller, E.\,W.\ Myers, D.\,J.\ Lipman:
\textit{Basic local alignment search tool.}
Journal of Molecular Biology, 215(3):403--410, 1990.

\bibitem{anderson1991infectious}
R.\,M.\ Anderson, R.\,M.\ May:
\textit{Infectious Diseases of Humans: Dynamics and Control.}
Oxford University Press, Oxford, 1991.

\bibitem{bennett1997strengths}
C.\,H.\ Bennett, E.\ Bernstein, G.\ Brassard, U.\ Vazirani:
\textit{Strengths and Weaknesses of Quantum Computing.}
SIAM Journal on Computing, 26(5):1510--1523, 1997.

\bibitem{blum1986simple}
L.\ Blum, M.\ Blum, M.\ Shub:
\textit{A Simple Unpredictable Pseudo-Random Number Generator.}
SIAM Journal on Computing, 15(2):364--383, 1986.

\bibitem{bloom1970space}
B.\,H.\ Bloom:
\textit{Space/time trade-offs in hash coding with allowable errors.}
Communications of the ACM, 13(7):422--426, 1970.

\bibitem{cohen2019certified}
J.\ Cohen, E.\ Rosenfeld, J.\,Z.\ Kolter:
\textit{Certified Adversarial Robustness via Randomized Smoothing.}
ICML 2019. Proceedings of Machine Learning Research, 97:1310--1320, 2019.

\bibitem{cook1971complexity}
S.\,A.\ Cook:
\textit{The Complexity of Theorem-Proving Procedures.}
Proc.\ 3rd Annual ACM Symposium on Theory of Computing (STOC), S.~151--158, 1971.

\bibitem{cormode2005improved}
G.\ Cormode, S.\ Muthukrishnan:
\textit{An Improved Data Stream Summary: The Count-Min Sketch and Its Applications.}
Journal of Algorithms, 55(1):58--75, 2005.

\bibitem{david1962games}
F.\,N.\ David:
\textit{Games, Gods and Gambling: The Origins and History of Probability and Statistical Ideas.}
Griffin, London, 1962.

\bibitem{dwork2014algorithmic}
C.\ Dwork, A.\ Roth:
\textit{The Algorithmic Foundations of Differential Privacy.}
Foundations and Trends in Theoretical Computer Science, 9(3--4):211--407, 2014.

\bibitem{farhi2014quantum}
E.\ Farhi, J.\ Goldstone, S.\ Gutmann:
\textit{A Quantum Approximate Optimization Algorithm.}
arXiv:1411.4028, 2014.

\bibitem{felsenstein1985confidence}
J.\ Felsenstein:
\textit{Confidence Limits on Phylogenies: An Approach Using the Bootstrap.}
Evolution, 39(4):783--791, 1985.

\bibitem{flajolet1985probabilistic}
P.\ Flajolet, G.\,N.\ Martin:
\textit{Probabilistic Counting Algorithms for Data Base Applications.}
Journal of Computer and System Sciences, 31(2):182--209, 1985.

\bibitem{goemans1995improved}
M.\,X.\ Goemans, D.\,P.\ Williamson:
\textit{Improved Approximation Algorithms for Maximum Cut and Satisfiability Problems Using Semidefinite Programming.}
Journal of the ACM, 42(6):1115--1145, 1995.

\bibitem{goldwasser1989knowledge}
S.\ Goldwasser, S.\ Micali, C.\ Rackoff:
\textit{The Knowledge Complexity of Interactive Proof Systems.}
SIAM Journal on Computing, 18(1):186--208, 1989.

\bibitem{grover1996fast}
L.\,K.\ Grover:
\textit{A fast quantum mechanical algorithm for database search.}
Proc.\ 28th ACM STOC, S.~212--219, 1996.

\bibitem{hastad2001some}
J.\ HÃ¥stad:
\textit{Some Optimal Inapproximability Results.}
Journal of the ACM, 48(4):798--859, 2001.

\bibitem{hastings1970monte}
W.\,K.\ Hastings:
\textit{Monte Carlo Sampling Methods Using Markov Chains and Their Applications.}
Biometrika, 57(1):97--109, 1970.

\bibitem{impagliazzo1997p}
R.\ Impagliazzo, A.\ Wigderson:
\textit{P = BPP if E Requires Exponential Circuits:
Derandomizing the XOR Lemma.}
Proc.\ 29th ACM STOC, S.~220--229, 1997.

\bibitem{johnson1974approximation}
D.\,S.\ Johnson:
\textit{Approximation Algorithms for Combinatorial Problems.}
Journal of Computer and System Sciences, 9(3):256--278, 1974.

\bibitem{karger1993global}
D.\,R.\ Karger:
\textit{Global Min-cuts in RNC, and Other Ramifications of a Simple Min-Cut Algorithm.}
Proc.\ 4th Annual ACM-SIAM Symposium on Discrete Algorithms (SODA), S.~21--30, 1993.

\bibitem{karger1996new}
D.\,R.\ Karger, C.\ Stein:
\textit{A New Approach to the Minimum Cut Problem.}
Journal of the ACM, 43(4):601--640, 1996.

\bibitem{kirkpatrick1983optimization}
S.\ Kirkpatrick, C.\,D.\ Gelatt, M.\,P.\ Vecchi:
\textit{Optimization by Simulated Annealing.}
Science, 220(4598):671--680, 1983.

\bibitem{kolmogorov1933grundbegriffe}
A.\,N.\ Kolmogorov:
\textit{Grundbegriffe der Wahrscheinlichkeitsrechnung.}
Springer, Berlin, 1933.
(Nachdruck: Springer, 1977.)

\bibitem{laplace1812theorie}
P.-S.\ Laplace:
\textit{Théorie analytique des probabilités.}
Courcier, Paris, 1812.
(3.~Auflage mit Supplements, 1820.)

\bibitem{mertens2006threshold}
S.\ Mertens, M.\ Mézard, R.\ Zecchina:
\textit{Threshold Values of Random K-SAT from the Cavity Method.}
Random Structures \& Algorithms, 28(3):340--373, 2006.

\bibitem{metropolis1953equation}
N.\ Metropolis, A.\,W.\ Rosenbluth, M.\,N.\ Rosenbluth, A.\,H.\ Teller, E.\ Teller:
\textit{Equation of State Calculations by Fast Computing Machines.}
Journal of Chemical Physics, 21(6):1087--1092, 1953.

\bibitem{miller1976riemann}
G.\,L.\ Miller:
\textit{Riemann's Hypothesis and Tests for Primality.}
Journal of Computer and System Sciences, 13(3):300--317, 1976.

\bibitem{moser2010constructive}
R.\,A.\ Moser, G.\ Tardos:
\textit{A Constructive Proof of the General Lovász Local Lemma.}
Journal of the ACM, 57(2):Article~11, 2010.

\bibitem{motwani1995randomized}
R.\ Motwani, P.\ Raghavan:
\textit{Randomized Algorithms.}
Cambridge University Press, Cambridge, 1995.

\bibitem{pugh1990skip}
W.\ Pugh:
\textit{Skip Lists: A Probabilistic Alternative to Balanced Trees.}
Communications of the ACM, 33(6):668--676, 1990.

\bibitem{rabin1980probabilistic}
M.\,O.\ Rabin:
\textit{Probabilistic Algorithm for Testing Primality.}
Journal of Number Theory, 12(1):128--138, 1980.

\bibitem{reingold2008undirected}
O.\ Reingold:
\textit{Undirected Connectivity in Log-Space.}
Journal of the ACM, 55(4):Article~17, 2008.

\bibitem{robbins1951stochastic}
H.\ Robbins, S.\ Monro:
\textit{A Stochastic Approximation Method.}
Annals of Mathematical Statistics, 22(3):400--407, 1951.

\bibitem{schwartz1980fast}
J.\,T.\ Schwartz:
\textit{Fast Probabilistic Algorithms for Verification of Polynomial Identities.}
Journal of the ACM, 27(4):701--717, 1980.

\bibitem{seidel1996randomized}
R.\ Seidel, C.\,R.\ Aragon:
\textit{Randomized Search Trees.}
Algorithmica, 16(4/5):464--497, 1996.

\bibitem{shannon1949communication}
C.\,E.\ Shannon:
\textit{Communication Theory of Secrecy Systems.}
Bell System Technical Journal, 28(4):656--715, 1949.

\bibitem{shor1994algorithms}
P.\,W.\ Shor:
\textit{Algorithms for Quantum Computation: Discrete Logarithms and Factoring.}
Proc.\ 35th Annual Symposium on Foundations of Computer Science (FOCS),
S.~124--134, 1994.

\bibitem{sipser2012introduction}
M.\ Sipser:
\textit{Introduction to the Theory of Computation.}
3.~Auflage.
Cengage Learning, Boston, 2012.

\bibitem{srivastava2014dropout}
N.\ Srivastava, G.\ Hinton, A.\ Krizhevsky, I.\ Sutskever, R.\ Salakhutdinov:
\textit{Dropout: A Simple Way to Prevent Neural Networks from Overfitting.}
Journal of Machine Learning Research, 15:1929--1958, 2014.

\bibitem{tashma2017explicit}
A.\ Ta-Shma:
\textit{Explicit, Almost Optimal, Epsilon-Balanced Codes.}
Proc.\ 49th ACM STOC, S.~238--251, 2017.

\bibitem{valiant1984theory}
L.\,G.\ Valiant:
\textit{A Theory of the Learnable.}
Communications of the ACM, 27(11):1134--1142, 1984.

\bibitem{williamson2011design}
D.\,P.\ Williamson, D.\,B.\ Shmoys:
\textit{The Design of Approximation Algorithms.}
Cambridge University Press, Cambridge, 2011.

\bibitem{zippel1979probabilistic}
R.\ Zippel:
\textit{Probabilistic Algorithms for Sparse Polynomials.}
Proc.\ International Symposium on Symbolic and Algebraic Computation (EUROSAM),
Lecture Notes in Computer Science 72, S.~216--226. Springer, 1979.

\end{thebibliography}

\end{document}

